%% ==========================================================================
%%  Chapter 9: Parallel Mechanisms, Constrained Dynamics, and Geometric Control
%% ==========================================================================
\chapter{Parallel Mechanisms, Constrained Dynamics, and Geometric Control}
\label{ch:parallel-mechanisms}

\begin{laymansbox}{Constraints as Geometric Scaffolds}{constraints_scaffolds}
  A golf club held in both hands, a parallel robot gripping an object, or legs pushing against the ground
  are all systems with \emph{constraints}: physical connections that remove degrees of freedom and
  couple subsystems. Rather than treating constraints as annoying algebraic equations to eliminate,
  we view them as \emph{geometric structures}—reductions of the configuration manifold, decompositions
  of force space, and sources of null-space redundancy. This chapter unifies the mathematical
  treatment of configuration manifolds, Lagrangian mechanics on submanifolds, and optimal control
  in the presence of constraints. The geometry reveals why and how complex motions can be decomposed,
  where redundancy lives, and how to design controls that respect the hidden structure.
\end{laymansbox}

\section{Configuration Manifolds and Holonomic Constraints}
\label{sec:ch9:config-manifolds}

\subsection{Unconstrained Configuration Space}

The starting point for any mechanical system is its \textbf{configuration manifold} $Q$, the space
of all kinematically possible positions of the system. For a system of $N$ rigid bodies in three
dimensions, the unconstrained configuration space is
\begin{equation}
  Q_{\text{free}} = (SE(3))^N,
\end{equation}
\index{Configuration manifold}%
where $SE(3)$ is the special Euclidean group (rigid body motions in 3D). Each copy of $SE(3)$
represents one body's position and orientation, encoded by a homogeneous transformation matrix
\begin{equation}
  T_i = \begin{bmatrix} R_i(\theta_i) & \mathbf{p}_i \\ 0 & 1 \end{bmatrix} \in SE(3),
\end{equation}
where $R_i \in SO(3)$ is the rotation matrix and $\mathbf{p}_i \in \mathbb{R}^3$ is the position.

The dimension of $Q_{\text{free}}$ is $\dim Q_{\text{free}} = 6N$, since each rigid body has 3
translational and 3 rotational degrees of freedom. For a system of $N$ links in a planar setting,
the dimension reduces to $\dim Q_{\text{free}} = 3N$.

\subsection{Holonomic and Nonholonomic Constraints}

\index{Constraint!holonomic}%
\index{Constraint!nonholonomic}%
A \textbf{holonomic constraint} is a geometric restriction on the configuration that can be expressed
as an algebraic equation:
\begin{equation}
  \phi_j(q) = 0, \quad j = 1, \ldots, c,
  \label{eq:ch9:holonomic-constraint}
\end{equation}
where $q \in Q_{\text{free}}$ and $\phi_j: Q_{\text{free}} \to \mathbb{R}$ is a smooth function. Examples include:
\begin{itemize}
  \item \textbf{Loop closure constraints}: In a multi-link chain or parallel mechanism, if two end-effectors
  must meet at a point (or achieve a fixed relative pose), this imposes holonomic constraints.
  \item \textbf{Contact constraints}: A foot in contact with the ground has its vertical position fixed
  and its horizontal position constrained (if friction prevents slipping).
  \item \textbf{Joint constraints}: Pin joints, sliding joints, and spherical joints all reduce the
  relative configuration space via holonomic constraints.
\end{itemize}

A \textbf{nonholonomic constraint}, by contrast, is a restriction on velocities that does not integrate
to a holonomic constraint:
\begin{equation}
  \sum_i a_i(q) \dot{q}_i = 0.
  \label{eq:ch9:nonholonomic-constraint}
\end{equation}
\index{Nonholonomic constraint}%
The canonical example is a rolling wheel or rolling foot: the velocity constraint ($\dot{x}_{\text{contact}} = v_{\text{wheel}}\dot{\theta}$)
does not prevent the wheel from moving to any configuration through some sequence of motions.
This chapter focuses on \textbf{holonomic constraints}; nonholonomic mechanics appear in Chapter~\ref{ch:nonholonomic-mechanics-if-present}
(if that chapter exists in future volumes).

\subsection{The Constraint Manifold as a Submanifold}

Let $\Phi: Q_{\text{free}} \to \mathbb{R}^c$ be the constraint map:
\begin{equation}
  \Phi(q) = \begin{bmatrix} \phi_1(q) \\ \vdots \\ \phi_c(q) \end{bmatrix}.
  \label{eq:ch9:constraint-map}
\end{equation}

The constraint manifold is the level set
\begin{equation}
  Q_c := \{q \in Q_{\text{free}} : \Phi(q) = 0\}.
  \label{eq:ch9:constraint-manifold}
\end{equation}

\index{Constraint Jacobian}%
The Jacobian of the constraint is the $c \times n$ matrix
\begin{equation}
  J_c(q) := \frac{\partial \Phi}{\partial q} \bigg|_q \in \mathbb{R}^{c \times n},
\end{equation}
where $n = \dim Q_{\text{free}}$.

\begin{proposition}[Regular Constraint Theorem]\label{prop:regular-constraint}
  Suppose that for all $q \in Q_c$, the Jacobian $J_c(q)$ has full row rank (rank $c$). Then:
  \begin{enumerate}
    \item $Q_c$ is a smooth submanifold of $Q_{\text{free}}$ of dimension $n_c := n - c$.
    \item The tangent space at any point $q \in Q_c$ is the kernel of $J_c(q)$:
    \begin{equation}
      T_q Q_c = \ker(J_c(q)) = \{\mathbf{v} \in T_q Q_{\text{free}} : J_c(q)\mathbf{v} = 0\}.
      \label{eq:ch9:tangent-constraint}
    \end{equation}
    \item The orthogonal complement (with respect to the kinetic-energy metric, see Section~\ref{sec:ch9:lagrangian})
    is spanned by the gradient vectors:
    \begin{equation}
      (T_q Q_c)^{\perp_M} = \text{span}\{\nabla_{q} \phi_1(q), \ldots, \nabla_{q} \phi_c(q)\}
      = \text{range}(J_c(q)^\top).
      \label{eq:ch9:orthogonal-complement}
    \end{equation}
  \end{enumerate}
\end{proposition}

\begin{proof}
  This is the regular level set theorem from differential topology. At a regular point $q \in Q_c$
  (where $J_c(q)$ has full row rank), the implicit function theorem guarantees that $Q_c$ is a manifold.
  The tangent space consists of tangent vectors $\mathbf{v}$ such that the directional derivative of
  each constraint vanishes: $\mathrm{d}\phi_j(\mathbf{v}) = J_c(q)\mathbf{v} = 0$. The orthogonal
  complement follows from the fact that the Euclidean metric (or kinetic-energy metric) makes the
  tangent space and its orthogonal complement complementary subspaces.
\end{proof}

\subsection{Constraint Singularities and Rank Loss}

When $J_c(q)$ loses rank at some point $q \in Q_c$, the constraint manifold becomes singular at that point.
Physically, this often corresponds to a \textbf{configuration singularity} or \textbf{grasp singularity}:
a posture where the mechanism loses rigidity or gains uncontrolled degrees of freedom.

For example, in a parallel robot gripper with 6 DOF and 6 constraints (position and orientation of the
end-effector), loss of rank in $J_c$ means the gripper can move in some direction while maintaining contact.
Such singularities are important to understand for control design: they represent configurations where
standard feedback cannot stabilize the system locally.

\subsection{Worked Example: Planar Two-Link Closed Chain}

\index{Closed-chain!two-link planar}%
Consider a simple planar closed-chain mechanism: two rigid links of lengths $\ell_1$ and $\ell_2$, with
joints at angles $\theta_1$ and $\theta_2$ (measured from the horizontal), constrained to meet at a common
endpoint. Let the base be at the origin.

The configuration space is $Q_{\text{free}} = [0, 2\pi) \times [0, 2\pi)$ (the pair $(\theta_1, \theta_2)$).
The constraint is that the endpoint of link 1 equals the endpoint of link 2, measured from a fixed reference:
\begin{align}
  \ell_1 \cos\theta_1 &= \ell_2 \cos\theta_2 + \Delta x, \\
  \ell_1 \sin\theta_1 &= \ell_2 \sin\theta_2 + \Delta y,
\end{align}
where $(\Delta x, \Delta y)$ is the fixed distance between the base points.

Rewrite as
\begin{align}
  \phi_1(\theta_1, \theta_2) &:= \ell_1 \cos\theta_1 - \ell_2 \cos\theta_2 - \Delta x = 0, \\
  \phi_2(\theta_1, \theta_2) &:= \ell_1 \sin\theta_1 - \ell_2 \sin\theta_2 - \Delta y = 0.
\end{align}

The constraint Jacobian is
\begin{equation}
  J_c = \begin{bmatrix}
    -\ell_1 \sin\theta_1 & \ell_2 \sin\theta_2 \\
    \ell_1 \cos\theta_1 & -\ell_2 \cos\theta_2
  \end{bmatrix}.
\end{equation}

For a regular point, we require $\text{rank}(J_c) = 2$. The determinant is
\begin{equation}
  \det(J_c) = \ell_1 \ell_2 (\sin\theta_1 \cos\theta_2 + \cos\theta_1 \sin\theta_2) = \ell_1 \ell_2 \sin(\theta_1 - \theta_2).
\end{equation}

The constraint manifold becomes singular when $\theta_1 = \theta_2$ (both links parallel) or $\theta_1 = \theta_2 + \pi$
(links anti-parallel). At these configurations, the two links are aligned and cannot be separated---the mechanism
has a locked or freely-moving singularity.

The constraint manifold $Q_c$ in the regular case is a 0-dimensional set (isolated points), but for special
geometries (e.g., $\ell_1 = \ell_2$), the manifold may be 1-dimensional.

\subsection{Revolute Joint Constraints on SE(3)$^2$}

\index{Joint constraint!revolute}%
A revolute joint constrains the relative motion of two bodies to rotation about a single axis. In the product
space $SE(3)^2$, where the first copy represents the frame of body 1 and the second represents body 2,
the joint imposes constraints on the relative transformation.

Let $g = g_1^{-1} g_2 \in SE(3)$ be the relative transformation. A revolute joint aligned with the $z$-axis
at a fixed point (the pivot) requires:
\begin{enumerate}
  \item \textbf{Position constraint}: The relative position of the pivot must be zero:
  \begin{equation}
    \phi_1(g) := p_{\text{rel}, x} = 0, \quad \phi_2(g) := p_{\text{rel}, y} = 0, \quad \phi_3(g) := p_{\text{rel}, z} = 0.
  \end{equation}
  These are 3 constraints on the translational part.

  \item \textbf{Rotational constraint}: The relative rotation must be about the $z$-axis only, so the rotation
  axis in the body frame must align with $z$. This is typically enforced as:
  \begin{equation}
    \phi_4(g) := R_{\text{rel}}(1, 2) = 0, \quad \phi_5(g) := R_{\text{rel}}(2, 1) = 0,
  \end{equation}
  where $R_{\text{rel}}(i, j)$ are the off-diagonal entries of the rotation part ensuring alignment.
\end{enumerate}

Thus, a single revolute joint imposes \textbf{5 constraints} on $SE(3)^2$ (3 translational + 2 rotational),
leaving \textbf{1 degree of freedom} (rotation about the axis).

\subsection{Holonomic vs. Nonholonomic Constraints: The Integrability Condition}

\index{Integrability condition}%
The distinction between holonomic and nonholonomic constraints lies in the concept of \textbf{integrability}.
A velocity constraint
\begin{equation}
  \sum_i a_i(q) \, \dot{q}_i = 0
\end{equation}
is called \textbf{holonomic} if there exists a function $\phi(q)$ such that
\begin{equation}
  a_i(q) = \frac{\partial \phi}{\partial q_i}.
\end{equation}

Equivalently, the 1-form $\alpha = \sum_i a_i(q) \, dq_i$ is \textbf{exact}, i.e., $\alpha = d\phi$ for some scalar function $\phi$.
By Frobenius's theorem, a 1-form is exact if and only if it is \textbf{closed}: $d\alpha = 0$.

For a single 1-form, this means
\begin{equation}
  \frac{\partial a_i}{\partial q_j} = \frac{\partial a_j}{\partial q_i} \quad \text{for all } i, j.
\end{equation}

\begin{example}[Holonomic Constraint: Grip Condition]
  Suppose a robot hand grips an object, and the contact point must remain at a fixed position in 3D space.
  The velocity constraint is $\dot{\mathbf{p}}_{\text{contact}} = 0$, which can be written as
  \begin{equation}
    \phi_x(q) := p_x(q) = 0, \quad \phi_y(q) := p_y(q) = 0, \quad \phi_z(q) := p_z(q) = 0.
  \end{equation}
  These are three holonomic constraints. The integrability condition is trivially satisfied because each
  constraint is a function of configuration only.
\end{example}

\begin{example}[Nonholonomic Constraint: Rolling Foot]
  A foot rolling on the ground without slipping satisfies $\mathbf{v}_{\text{contact}} = 0$ (no slipping).
  But the foot can also \emph{roll} (angular motion about the contact line). The velocity constraint is
  $v_x(q, \dot{q}) = \dot{q}_{\text{roll}} \times r = 0$ (approximately), which does not integrate to
  a holonomic constraint---the foot can reach different positions through different rolling paths.

  The 1-form is not exact: $d\alpha \ne 0$, and the constraint manifold is not a simple level set.
\end{example}

For the golf swing and grasp mechanics studied in this chapter, we focus on \textbf{holonomic constraints}
(grip points, contact patches, joint limits), while acknowledging that rolling contacts and sliding are
inherently nonholonomic and require specialized treatment.

\section{Lagrangian Mechanics on Constrained Manifolds}
\label{sec:ch9:lagrangian}

\subsection{Kinetic Energy and the Riemannian Metric}

On the unconstrained tangent bundle $T Q_{\text{free}}$, the kinetic energy is
\begin{equation}
  T = \frac{1}{2} \dot{q}^\top M(q) \dot{q},
  \label{eq:ch9:kinetic-energy}
\end{equation}
\index{Mass matrix}%
where $M(q) \in \mathbb{R}^{n \times n}$ is the symmetric, positive-definite \textbf{mass matrix}
(or \textbf{inertia matrix}). For a chain of rigid bodies, the entries of $M$ depend only on the
configuration $q$:
\begin{equation}
  M_{ij}(q) = \sum_{k} \left( m_k J_{k,i}^\top J_{k,j} + \text{tr}(R_k R_k^\top I_k J_{\omega,i}^\top J_{\omega,j}) \right),
\end{equation}
where $m_k$ is the mass of link $k$, $J_{k,i}$ is the translational Jacobian with respect to $q_i$,
$I_k$ is the inertia tensor in body frame, $R_k$ is the current rotation, and $J_{\omega,i}$ is the
rotational Jacobian.

The mass matrix induces a \textbf{Riemannian metric} on $Q_{\text{free}}$:
\begin{equation}
  \langle \mathbf{u}, \mathbf{v} \rangle_M := \mathbf{u}^\top M(q) \mathbf{v} \quad \text{for } \mathbf{u}, \mathbf{v} \in T_q Q.
  \label{eq:ch9:metric}
\end{equation}
This metric encodes the cost of acceleration: moving fast in a direction requires more kinetic energy
if $M$ is large in that direction.

\subsection{Euler-Lagrange Equations with Constraints}

The \textbf{Lagrangian} is $L = T - V$, where $V(q)$ is the potential energy (usually gravitational).
On the unconstrained space, the Euler-Lagrange equations are
\begin{equation}
  \frac{d}{dt}\frac{\partial L}{\partial \dot{q}} - \frac{\partial L}{\partial q} = \mathbf{F}_{\text{ext}} + \mathbf{F}_{\text{control}},
\end{equation}
where $\mathbf{F}_{\text{ext}}$ and $\mathbf{F}_{\text{control}}$ are external and control forces.

When constraints $\Phi(q) = 0$ are active, we introduce Lagrange multipliers $\lambda \in \mathbb{R}^c$
and write the constrained Euler-Lagrange equations:
\begin{equation}
  M(q)\ddot{q} + C(q,\dot{q})\dot{q} + \nabla V(q) = J_c(q)^\top \lambda + B(q)\mathbf{u},
  \label{eq:ch9:euler-lagrange-constrained}
\end{equation}
\index{Lagrange multipliers}%
where:
\begin{itemize}
  \item $C(q,\dot{q})$ is the Coriolis-centrifugal matrix, computed from the mass matrix via
  \begin{equation}
    C(q,\dot{q}) = \frac{\partial M}{\partial t}\bigg|_{\text{convected}} = \sum_i \dot{q}_i \frac{\partial M}{\partial q_i}.
    \label{eq:ch9:coriolis}
  \end{equation}
  \item $\nabla V(q)$ is the gravitational gradient.
  \item $J_c(q)^\top \lambda$ is the constraint force (normal to the constraint manifold).
  \item $B(q) \in \mathbb{R}^{n \times m}$ is the input matrix, and $\mathbf{u} \in \mathbb{R}^m$ are
  the control inputs (torques at actuated joints).
\end{itemize}

The constraint forces $J_c(q)^\top \lambda$ enforce the acceleration constraint:
\begin{equation}
  J_c(q) \ddot{q} + \frac{d}{dt}(J_c(q))\dot{q} = 0,
  \label{eq:ch9:accel-constraint}
\end{equation}
which follows by differentiating $\Phi(q(t)) = 0$ twice in time.

\subsection{Christoffel Symbols and Geometric Meaning of Coriolis}

\index{Christoffel symbol}%
On a Riemannian manifold, the \textbf{Christoffel symbols} $\Gamma^k_{ij}$ describe how the basis vectors
of the tangent space change as we move through the manifold. In coordinates, they are given by
\begin{equation}
  \Gamma^k_{ij} = \frac{1}{2} M^{kl} \left( \frac{\partial M_{lj}}{\partial q_i} + \frac{\partial M_{li}}{\partial q_j}
  - \frac{\partial M_{ij}}{\partial q_l} \right),
  \label{eq:ch9:christoffel}
\end{equation}
where $M^{kl}$ is the $(k,l)$-entry of $M^{-1}$.

The Coriolis-centrifugal acceleration in configuration space can be expressed in terms of Christoffel symbols:
\begin{equation}
  (C(q,\dot{q})\dot{q})_i = -\sum_{j,k} \Gamma^j_{ik}\dot{q}_j\dot{q}_k.
  \label{eq:ch9:coriolis-christoffel}
\end{equation}

\begin{intuitionbox}
  The Coriolis and centrifugal terms are \emph{not} ``real'' forces acting on the system. Rather,
  they are artifacts of our choice of coordinates on the curved configuration manifold. They arise
  because the basis vectors $\partial/\partial q_i$ are not constant; they rotate and stretch as
  we move through $Q$. A coordinate-free description would eliminate them, but working in local
  coordinates makes them explicit and important for numerical computation.
\end{intuitionbox}

\subsection{Christoffel Symbol Computation: A 2-DOF Example}

\index{Christoffel symbol!example}%
Consider a 2-DOF system with generalized coordinates $(q_1, q_2)$ and mass matrix
\begin{equation}
  M(q) = \begin{bmatrix} m_1 + m_2 & m_2 \cos(q_2) \\ m_2 \cos(q_2) & m_2 \end{bmatrix},
\end{equation}
where $m_1$ and $m_2$ are link masses, and the $(1,2)$ entry includes a position-dependent coupling.

The Christoffel symbols \eqref{eq:ch9:christoffel} require computing the inverse:
\begin{equation}
  M^{-1} = \frac{1}{\det(M)} \begin{bmatrix} m_2 & -m_2 \cos(q_2) \\ -m_2 \cos(q_2) & m_1 + m_2 \end{bmatrix},
\end{equation}
where $\det(M) = m_1 m_2 + m_2^2(1 - \cos^2(q_2)) = m_1 m_2 + m_2^2 \sin^2(q_2)$.

For the term $\Gamma^1_{22}$ (which produces centrifugal acceleration in direction 1 due to velocity in direction 2):
\begin{align}
  \Gamma^1_{22} &= \frac{1}{2} (M^{-1})_{1l} \left( \frac{\partial M_{l2}}{\partial q_1} + \frac{\partial M_{l2}}{\partial q_1}
  - \frac{\partial M_{22}}{\partial q_l} \right), \\
  &= \frac{1}{2} (M^{-1})_{11} \left( 0 + 0 - 0 \right) + \frac{1}{2} (M^{-1})_{12} \left( 0 + 0 - 0 \right), \\
  &= 0.
\end{align}

But $\Gamma^1_{12}$ (Coriolis coupling):
\begin{align}
  \Gamma^1_{12} &= \frac{1}{2} (M^{-1})_{1l} \left( \frac{\partial M_{l2}}{\partial q_1} + \frac{\partial M_{l1}}{\partial q_2}
  - \frac{\partial M_{12}}{\partial q_l} \right), \\
  &= \frac{1}{2} (M^{-1})_{12} \left( \frac{\partial M_{22}}{\partial q_1} + 0 - \frac{\partial M_{12}}{\partial q_1} \right).
\end{align}

Since $\frac{\partial M_{22}}{\partial q_1} = 0$ and $\frac{\partial M_{12}}{\partial q_1} = -m_2 \sin(q_2)$:
\begin{equation}
  \Gamma^1_{12} = \frac{1}{2} (M^{-1})_{12} \cdot m_2 \sin(q_2) = \frac{m_2 \sin(q_2)}{2 \det(M)} \cdot m_2 \sin(q_2).
\end{equation}

This produces the Coriolis term $(C(q,\dot{q})\dot{q})_1 = -\Gamma^1_{12} \dot{q}_1 \dot{q}_2$ and centrifugal terms,
which are essential for accurate simulation and control.

\subsection{Routh Reduction for Cyclic Coordinates}

\index{Routh reduction}%
\index{Cyclic coordinate}%
When the Lagrangian $L(q, \dot{q}) = T(q, \dot{q}) - V(q)$ is invariant under a symmetry group action,
certain generalized coordinates (called \textbf{cyclic coordinates}) do not appear explicitly in $L$.

For such a coordinate $q_{\text{cyc}}$, the Euler-Lagrange equation gives
\begin{equation}
  \frac{d}{dt}\frac{\partial L}{\partial \dot{q}_{\text{cyc}}} = 0,
\end{equation}
so the generalized momentum is conserved:
\begin{equation}
  p_{\text{cyc}} = \frac{\partial L}{\partial \dot{q}_{\text{cyc}}} = \text{constant}.
\end{equation}

The \textbf{Routhian} is constructed by eliminating the cyclic velocity:
\begin{equation}
  R(q_{\text{shape}}, \dot{q}_{\text{shape}}, p_{\text{cyc}}) := \dot{q}_{\text{cyc}} p_{\text{cyc}} - L(q_{\text{shape}}, \dot{q}_{\text{shape}}, \dot{q}_{\text{cyc}}(p_{\text{cyc}})),
\end{equation}
where $\dot{q}_{\text{cyc}} = \dot{q}_{\text{cyc}}(q_{\text{shape}}, \dot{q}_{\text{shape}}, p_{\text{cyc}})$ is determined from the constraint.

The dynamics of the shape variables reduce to
\begin{equation}
  \frac{d}{dt}\frac{\partial R}{\partial \dot{q}_{\text{shape}}} - \frac{\partial R}{\partial q_{\text{shape}}} = 0,
\end{equation}
a system of lower dimension. This is crucial for understanding the golf swing: the total angular momentum
is conserved (cyclic coordinate), and the swing dynamics can be reduced to the shape changes of the arms,
with the rotation of the torso following deterministically.

\subsection{Index-3 DAE Formulation and Stabilization}

\index{DAE!index-3}%
The system of equations
\begin{align}
  M(q)\ddot{q} + C(q,\dot{q})\dot{q} + \nabla V(q) &= J_c(q)^\top \lambda + B(q)\mathbf{u}, \label{eq:ch9:dae-1}\\
  \Phi(q) &= 0, \label{eq:ch9:dae-2}
\end{align}
is a \textbf{differential-algebraic equation (DAE)} of \textbf{index 3}. The index is the minimum
number of times the algebraic constraint must be differentiated to derive an ODE for $\lambda$.

To understand why, differentiate the constraint \eqref{eq:ch9:dae-2} in time:
\begin{equation}
  \dot{\Phi} = J_c(q) \dot{q} = 0. \label{eq:ch9:dae-vel}
\end{equation}

Differentiate again:
\begin{equation}
  \ddot{\Phi} = J_c(q) \ddot{q} + \dot{J}_c(q,\dot{q}) \dot{q} = 0. \label{eq:ch9:dae-acc}
\end{equation}

Now substitute \eqref{eq:ch9:dae-1} into \eqref{eq:ch9:dae-acc}:
\begin{equation}
  J_c(q) M(q)^{-1} \left( J_c(q)^\top \lambda + B(q)\mathbf{u} - C(q,\dot{q})\dot{q} - \nabla V(q) \right) + \dot{J}_c(q,\dot{q}) \dot{q} = 0.
\end{equation}

This is an equation for $\lambda$. Thus, we need three differentiations (position $\to$ velocity $\to$ acceleration)
to obtain an ODE system. The system is of \textbf{index 3}.

\subsubsection{Baumgarte Stabilization}

To solve this numerically, the \textbf{Baumgarte stabilization} method augments the acceleration constraint
with corrective terms proportional to the position and velocity constraint violations:
\begin{equation}
  J_c(q)\ddot{q} + \dot{J}_c(q,\dot{q})\dot{q} + 2\alpha \dot{\Phi}(q,\dot{q}) + \beta^2 \Phi(q) = 0,
  \label{eq:ch9:baumgarte}
\end{equation}
where $\alpha, \beta > 0$ are tuning parameters (typically $\alpha = \beta = 1$ or chosen to damp oscillations).

The correction terms $2\alpha \dot{\Phi} + \beta^2 \Phi$ can be interpreted as a feedback law: if the constraint
is violated, the correction grows to drive the state back onto the manifold. This reduces the effective index
to 1 and suppresses exponential growth of constraint violations due to numerical integration error.

The modified constraint can be rewritten as a second-order ODE in $\Phi$:
\begin{equation}
  \ddot{\Phi} + 2\alpha \dot{\Phi} + \beta^2 \Phi = 0,
\end{equation}
which is stable (like a damped harmonic oscillator) if $\alpha, \beta > 0$.

\subsubsection{Projection Method}

Alternatively, the \textbf{projection method} projects the unconstrained accelerations onto the constraint
manifold at each time step:
\begin{equation}
  \ddot{q}_{\text{proj}} = \left( I - J_c(q)^\top (J_c(q) M^{-1}(q) J_c(q)^\top)^{-1} J_c(q) M^{-1}(q) \right) \ddot{q}_{\text{unconstrained}},
\end{equation}
ensuring exact constraint satisfaction up to machine precision. This is more accurate but computationally
expensive (requires solving a linear system at each time step).

\section{Control-Affine Systems on Configuration Manifolds}
\label{sec:ch9:control-affine}

\subsection{Control-Affine Dynamics and Drift Field}

\index{Control-affine system}%
A system is \textbf{control-affine} if the dynamics can be written as
\begin{equation}
  \ddot{q} = f(q,\dot{q}) + \sum_{i=1}^{m} g_i(q) u_i,
  \label{eq:ch9:control-affine}
\end{equation}
where $f: TQ \to TQ$ is the \textbf{drift} vector field (dynamics with zero control) and $g_i: Q \to \mathbb{R}^n$
are the \textbf{input fields}.

From the constrained Euler-Lagrange equation \eqref{eq:ch9:euler-lagrange-constrained}, we identify:
\begin{align}
  f(q,\dot{q}) &= -M(q)^{-1} \left( C(q,\dot{q})\dot{q} + \nabla V(q) - J_c(q)^\top \lambda \right), \label{eq:ch9:drift}\\
  g_i(q) &= M(q)^{-1} B(q)e_i, \label{eq:ch9:input-field}
\end{align}
where $e_i$ is the $i$-th standard basis vector.

Note that the constraint forces $J_c(q)^\top \lambda$ are part of the drift: they are not ``controlled''
in the usual sense but rather reaction forces that enforce the constraint. These must be computed from
the constraint equation \eqref{eq:ch9:accel-constraint}.

\subsection{Lie Bracket and Controllability}

\index{Lie bracket}%
On the tangent bundle $TQ$, we define vector fields. The \textbf{Lie bracket} of two vector fields $X$ and $Y$
measures the difference between flowing along $X$ then $Y$ versus $Y$ then $X$:
\begin{equation}
  [X, Y] = X \circ Y - Y \circ X.
\end{equation}

In local coordinates, if $X = \sum_i X_i \frac{\partial}{\partial q_i}$ and $Y = \sum_j Y_j \frac{\partial}{\partial q_j}$, then
\begin{equation}
  [X, Y] = \sum_{i,j} \left( X_i \frac{\partial Y_j}{\partial q_i} - Y_i \frac{\partial X_j}{\partial q_i} \right) \frac{\partial}{\partial q_j}.
  \label{eq:ch9:lie-bracket-coords}
\end{equation}

\begin{example}[Explicit Lie Bracket Computation]
  Let $X = [x_2; -x_1]$ and $Y = [0; 1]$ (as column vectors). Compute $[X, Y]$.

  In coordinates on $\mathbb{R}^2$, we have
  \begin{align}
    X &= x_2 \frac{\partial}{\partial x_1} - x_1 \frac{\partial}{\partial x_2}, \\
    Y &= \frac{\partial}{\partial x_2}.
  \end{align}

  Applying $X$ to the components of $Y$:
  \begin{align}
    X(Y_1) &= X(0) = 0, \\
    X(Y_2) &= X(1) = 0.
  \end{align}

  Applying $Y$ to the components of $X$:
  \begin{align}
    Y(X_1) &= Y(x_2) = \frac{\partial x_2}{\partial x_2} = 1, \\
    Y(X_2) &= Y(-x_1) = \frac{\partial(-x_1)}{\partial x_2} = 0.
  \end{align}

  Using \eqref{eq:ch9:lie-bracket-coords}:
  \begin{equation}
    [X, Y] = \left( 0 - 1 \right) \frac{\partial}{\partial x_1} + \left( 0 - 0 \right) \frac{\partial}{\partial x_2}
    = -\frac{\partial}{\partial x_1} = [1; 0].
  \end{equation}

  This new direction is independent of $X$ and $Y$, showing that two inputs can generate motion in a third direction
  via iterated Lie brackets.
\end{example}

For the system \eqref{eq:ch9:control-affine}, the controllability condition (Hermann-Nagano, Chow's theorem)
states that the system is (generically) controllable if the Lie algebra generated by $\{g_1, \ldots, g_m\}$
and their iterated Lie brackets spans the entire tangent space at almost every point:
\begin{equation}
  \text{span}\left\{g_i, [g_i, g_j], [[g_i,g_j], g_k], \ldots\right\}\bigg|_q = T_q(TQ).
  \label{eq:ch9:controllability-condition}
\end{equation}

\begin{example}[Non-Holonomic Car on a Plane]
  A car with front-wheel steering has state $(x, y, \theta, \phi)$ where $(x,y)$ is position,
  $\theta$ is heading, and $\phi$ is front-wheel angle. The kinematics are
  \begin{equation}
    \begin{bmatrix} \dot{x} \\ \dot{y} \\ \dot{\theta} \\ \dot{\phi} \end{bmatrix}
    = \begin{bmatrix} \cos\theta \\ \sin\theta \\ \tan\phi/L \\ 0 \end{bmatrix} v + \begin{bmatrix} 0 \\ 0 \\ 0 \\ 1 \end{bmatrix} \omega,
  \end{equation}
  where $v$ is the forward speed and $\omega$ is the steering rate.

  The drift field is zero: $f = 0$. The input fields are $g_1 = (\cos\theta, \sin\theta, \tan\phi/L, 0)^\top$
  and $g_2 = (0, 0, 0, 1)^\top$.

  Computing Lie brackets: $[g_2, g_1]$ contains components that allow changes in $x$ and $y$ while
  maintaining zero velocity (by oscillating steering), demonstrating that the car is controllable despite
  having only 2 inputs and 4 degrees of freedom.
\end{example}

\subsection{Geodesic Spray and the Levi-Civita Connection}

\index{Geodesic spray}%
\index{Levi-Civita connection}%
The drift field $f(q, \dot{q}) = -M(q)^{-1}(C(q,\dot{q})\dot{q} + \nabla V(q))$ on the tangent bundle $TQ$
has a geometric interpretation via the Levi-Civita connection.

The \textbf{geodesic spray} associated with the Riemannian metric $g = M$ is the second-order ODE on $TQ$ defined by
\begin{equation}
  \ddot{q}^k + \Gamma^k_{ij}(q) \dot{q}^i \dot{q}^j = 0,
  \label{eq:ch9:geodesic-spray}
\end{equation}
where $\Gamma^k_{ij}$ are the Christoffel symbols. This describes motion along geodesics of the metric.

In our constrained system with gravity, the equation of motion becomes
\begin{equation}
  \ddot{q}^k + \Gamma^k_{ij}(q) \dot{q}^i \dot{q}^j = -M^{-1}(q) \nabla V(q) + M^{-1}(q) B(q) u.
\end{equation}

The drift field $f(q, \dot{q})$ encapsulates both the Christoffel symbols (geodesic curvature) and the gradient
of the potential. This connection reveals that the Coriolis and centrifugal terms, though appearing as ``fictitious forces''
in a coordinate-dependent view, are actually manifestations of the curvature of the configuration manifold with respect
to the kinetic-energy metric \cite{BulloLewis2004}.

\subsection{Input Distribution and Metric Orthogonal Complement}

\index{Input distribution}%
The \textbf{input distribution} is the span of the input fields:
\begin{equation}
  \Delta_{\text{input}}(q) := \text{span}\{g_1(q), \ldots, g_m(q)\}.
\end{equation}

The orthogonal complement with respect to the metric $M(q)$ is
\begin{equation}
  \Delta_{\text{input}}^{\perp_M}(q) := \{\mathbf{v} \in T_q Q : \mathbf{v}^\top M(q) g_i(q) = 0 \text{ for all } i = 1, \ldots, m\}.
\end{equation}

Physically, directions in $\Delta_{\text{input}}^{\perp_M}$ are ``expensive to actuate''---they require large inputs
relative to the kinetic energy gained. In redundant manipulators, null-space motions are chosen to avoid high-cost
directions.

\subsection{Passivity and Energy Storage}

\index{Passivity}%
A mechanical system is \textbf{passive} if the power supplied by external forces equals the rate of change of
stored energy (kinetic + potential):
\begin{equation}
  u^\top \tau_{\text{ext}} = \frac{d}{dt}(T + V),
  \label{eq:ch9:passivity}
\end{equation}
where $u$ are control inputs (torques or forces) and $\tau_{\text{ext}}$ is the external wrench.

Integrating over time:
\begin{equation}
  \int_0^t u(\sigma)^\top \tau_{\text{ext}}(\sigma) \, d\sigma = (T(t) + V(t)) - (T(0) + V(0)).
\end{equation}

For a golf swing, the golfer supplies energy through muscle activation. The swing is passive if the mechanical
energy at ball contact equals the input energy minus losses (air resistance, grip friction). Energy-optimal swings
minimize the input work subject to achieving a target ball velocity.

Passivity-based control exploits this structure: by designing feedback laws that preserve or reduce total energy,
we ensure stability without detailed knowledge of system parameters \cite{BulloLewis2004}.

\section{Null Spaces, Redundancy, and Families of Solutions}
\label{sec:ch9:null-space}

\subsection{Algebraic Redundancy and the Nullspace of the Input Matrix}

\index{Redundancy!algebraic}%
When the number of inputs $m$ is less than the dimension of the configuration space $n_c$ (the dimension
of the constrained manifold $Q_c$), the system is \textbf{redundant}: there are multiple ways to
achieve the same acceleration.

More precisely, suppose we seek to track a desired acceleration $\ddot{q}_d \in T Q_c$. The control law
\begin{equation}
  B(q) u = M(q)(\ddot{q}_d - f(q,\dot{q}))
  \label{eq:ch9:control-law}
\end{equation}
may have no solution if the right-hand side is not in the range of $B(q)$. But if $m < n_c$, there
are typically infinite solutions (a family parameterized by the null space):
\begin{equation}
  u = B(q)^+ M(q)(\ddot{q}_d - f(q,\dot{q})) + (I - B(q)^+ B(q)) z,
  \label{eq:ch9:general-control}
\end{equation}
where $B(q)^+$ is the pseudoinverse and $z$ is an arbitrary vector in $\mathbb{R}^m$.

\subsection{Velocity Constraints and the Constraint Jacobian}

The constraint equation $\Phi(q) = 0$ implies, after differentiation,
\begin{equation}
  J_c(q) \dot{q} = 0.
  \label{eq:ch9:velocity-constraint}
\end{equation}

This means any admissible velocity must lie in the null space of $J_c(q)$:
\begin{equation}
  \dot{q} \in \ker(J_c(q)).
\end{equation}

The orthogonal complement (with respect to metric $M$) is the range of $J_c(q)^\top$:
\begin{equation}
  (\ker(J_c(q)))^{\perp_M} = \text{range}(J_c(q)^\top).
\end{equation}

This decomposition is fundamental: any force can be uniquely written as
\begin{equation}
  \mathbf{F} = \mathbf{F}_{\text{useful}} + \mathbf{F}_{\text{reaction}},
\end{equation}
where $\mathbf{F}_{\text{useful}} \in \ker(J_c(q))^\top = \text{range}(J_c(q))$ acts along the constraint
and $\mathbf{F}_{\text{reaction}} \in \text{range}(J_c(q)^\top)$ is normal to the constraint manifold.

\subsection{Pseudoinverse and Minimum-Norm Solutions}

\index{Pseudoinverse}%
Given a desired velocity $\mathbf{v}_d$ on the constraint manifold $Q_c$, we may seek a velocity in the full
unconstrained space that approximates $\mathbf{v}_d$. The minimum-norm solution is
\begin{equation}
  \dot{q}_{\text{general}} = J_c(q)^+ \mathbf{v}_d + (I - J_c(q)^+ J_c(q)) \mathbf{z},
  \label{eq:ch9:min-norm-velocity}
\end{equation}
where $J_c(q)^+ = J_c(q)^\top (J_c(q)J_c(q)^\top)^{-1}$ is the Moore-Penrose pseudoinverse and
$\mathbf{z} \in \mathbb{R}^{n_c}$ is an arbitrary vector (selecting which direction in the null space).

The first term $J_c(q)^+ \mathbf{v}_d$ is the \textbf{minimum-norm} solution: it requires the least
kinetic energy (under metric $M$) to achieve the desired output. The null-space component $(I - J_c(q)^+ J_c(q))\mathbf{z}$
is \textbf{motor abundance}: motions in this subspace do not change the output but consume energy and
can be optimized separately.

\subsection{Wrench Space and Constraint Wrench Space}

\index{Wrench space}%
\index{Constraint wrench space}%
In the dual space (forces and moments), a \textbf{wrench} is an element of the cotangent space $T_q^* Q$.
A wrench can be written as $\mathbf{W} = \lambda \in \mathbb{R}^m$, where $\mathbf{W}^\top \mathbf{v}$ represents
the power (work per unit time) exerted by the wrench in direction $\mathbf{v}$.

The \textbf{wrench space} (all possible wrenches) is $\mathcal{W}_{\text{all}} = T_q^* Q \cong \mathbb{R}^m$.

The \textbf{constraint wrench space} is the set of wrenches that do no work on velocities in the constraint manifold:
\begin{equation}
  \mathcal{W}_c = \{\mathbf{W} \in T_q^* Q : \mathbf{W}^\top \mathbf{v} = 0 \text{ for all } \mathbf{v} \in T_q Q_c\}
  = \text{range}(J_c(q)^\top).
  \label{eq:ch9:constraint-wrench-space}
\end{equation}

These are the constraint reaction forces: forces normal to the constraint manifold that enforce the constraints.

The dimension formula:
\begin{equation}
  \dim(\mathcal{W}_c) = \text{rank}(J_c(q)^\top) = \text{rank}(J_c(q)) = c.
\end{equation}

Conversely, the \textbf{feasible wrench space} (wrenches that can do work on the constraint manifold) is the
orthogonal complement:
\begin{equation}
  \mathcal{W}_{\text{feas}} = (\mathcal{W}_c)^{\perp} = \ker(J_c(q)^\top),
\end{equation}
with dimension $\dim(\mathcal{W}_{\text{feas}}) = m - \text{rank}(J_c) = m - c = n_c$ (the dimension of the constraint manifold).

\subsubsection{Minimum Muscle Stress Criterion}

\index{Crowninshield-Brand criterion}%
In biomechanics, the inverse problem is to determine muscle forces given kinematics. With $N$ muscles and fewer
constraints (from joints), the problem is redundant: infinitely many muscle force combinations produce the same
joint torques.

The \textbf{Crowninshield \& Brand (1981)} criterion minimizes muscle stress (force per physiological cross-sectional area):
\begin{equation}
  \min_{f_m} \sum_i \frac{f_{m,i}}{A_i}^p,
\end{equation}
subject to
\begin{equation}
  \sum_i R_i f_{m,i} = \tau_j,
\end{equation}
where $f_{m,i}$ is the force in muscle $i$, $A_i$ is its cross-sectional area, $R_i$ is its moment arm, and
$\tau_j$ are joint torques. The exponent $p$ (typically $p = 2$) emphasizes distributing load evenly across muscles.

This is a convex quadratic program (if $p = 2$) that can be solved efficiently. The solution lies in the null space
of the moment-arm matrix and represents a physiologically plausible distribution of muscle activations.

\subsubsection{Optimal Null-Space Velocity for Torque Minimization}

\index{Null-space optimization}%
In a redundant manipulator with $m > n_c$ actuators, we seek to achieve a desired acceleration $\ddot{q}_d$
while minimizing control effort (torque squared).

The general solution is
\begin{equation}
  u = B^\dagger M(\ddot{q}_d - f) + (I - B^\dagger B) z,
  \label{eq:ch9:control-general-redundant}
\end{equation}
where $B^\dagger$ is the pseudoinverse and $z \in \mathbb{R}^m$ is arbitrary (null-space component).

To minimize $\|u\|^2 = \|B^\dagger M(\ddot{q}_d - f)\|^2 + \|(I - B^\dagger B) z\|^2$, we choose
\begin{equation}
  z^* = 0,
\end{equation}
giving the minimum-norm control.

But we may also want to minimize joint torques or muscle activations, which is reflected in a weighted cost:
\begin{equation}
  \|u\|_W^2 = u^\top W u,
\end{equation}
where $W$ is a positive-definite weight matrix.

The optimal null-space velocity is then
\begin{equation}
  z^* = -(I - B^\dagger B) \left( \frac{dW}{dq} \text{vec}(B^\dagger M(\ddot{q}_d - f)) \right),
  \label{eq:ch9:optimal-null-space}
\end{equation}
a corrective term that minimizes the predicted cost gradient. This enables motion planning that optimizes
over redundancy online.

\section{Screw Theory, Lie Groups, and Rigid Body Motion}
\label{sec:ch9:screw-theory}

\subsection{The Lie Algebra se(3) and Twists}

\index{Lie algebra!se(3)}%
\index{Twist!rigid body}%
The Lie algebra of $SE(3)$ is $\mathfrak{se}(3)$, the space of infinitesimal rigid motions (twists).
A twist is a pair $\xi = (\boldsymbol{\omega}, \mathbf{v}) \in \mathbb{R}^3 \times \mathbb{R}^3$, where
$\boldsymbol{\omega}$ is angular velocity and $\mathbf{v}$ is linear velocity.

In matrix form, a twist is represented as
\begin{equation}
  \hat{\xi} = \begin{bmatrix} [\boldsymbol{\omega}]_\times & \mathbf{v} \\ 0 & 0 \end{bmatrix} \in \mathbb{R}^{4 \times 4},
  \label{eq:ch9:twist-matrix}
\end{equation}
where $[\boldsymbol{\omega}]_\times$ is the skew-symmetric matrix
\begin{equation}
  [\boldsymbol{\omega}]_\times = \begin{bmatrix} 0 & -\omega_3 & \omega_2 \\ \omega_3 & 0 & -\omega_1 \\ -\omega_2 & \omega_1 & 0 \end{bmatrix}.
\end{equation}

\subsection{The Exponential Map: From Algebra to Group}

\index{Exponential map!SE(3)}%
The exponential map $\exp: \mathfrak{se}(3) \to SE(3)$ takes a twist $\xi$ and produces a rigid transformation:
\begin{equation}
  \exp(\hat{\xi}) = \begin{bmatrix} R & \mathbf{p} \\ 0 & 1 \end{bmatrix},
  \label{eq:ch9:exp-se3}
\end{equation}
where the rotation $R$ and translation $\mathbf{p}$ are computed using Rodrigues' formula:
\begin{align}
  R &= I + \sin(\|\boldsymbol{\omega}\|)\frac{[\boldsymbol{\omega}]_\times}{\|\boldsymbol{\omega}\|}
  + (1-\cos(\|\boldsymbol{\omega}\|))\frac{[\boldsymbol{\omega}]_\times^2}{\|\boldsymbol{\omega}\|^2},\\
  \mathbf{p} &= \left( I + (1-\cos(\|\boldsymbol{\omega}\|))\frac{[\boldsymbol{\omega}]_\times}{\|\boldsymbol{\omega}\|^2}
  + (\|\boldsymbol{\omega}\| - \sin(\|\boldsymbol{\omega}\|))\frac{[\boldsymbol{\omega}]_\times^2}{\|\boldsymbol{\omega}\|^3} \right)\mathbf{v}.
\end{align}

\subsection{Adjoint and Coadjoint Actions}

\index{Adjoint action}%
\index{Coadjoint action}%
When a rigid transformation $g \in SE(3)$ acts on a twist $\xi \in \mathfrak{se}(3)$, the result is
a \textbf{transformed twist}:
\begin{equation}
  \text{Ad}_g(\xi) = g \hat{\xi} g^{-1},
\end{equation}
which in matrix form gives a new twist in the body frame.

Explicitly, if $g = [R \, \mathbf{p}; 0 \, 1] \in SE(3)$ and $\hat{\xi} = [[\omega]_\times \, \mathbf{v}; 0 \, 0]$, then
\begin{equation}
  \text{Ad}_g(\hat{\xi}) = \begin{bmatrix} R[\omega]_\times R^\top & R\mathbf{v} \\ 0 & 0 \end{bmatrix},
  \label{eq:ch9:adjoint-matrix}
\end{equation}
where the rotated angular velocity is $\omega' = R\omega$ and the transformed linear velocity is $\mathbf{v}' = R\mathbf{v} - [\mathbf{p}]_\times R\omega$.

The \textbf{coadjoint action} acts on wrench space (dual to twists):
\begin{equation}
  \text{Ad}^*_g(\mathbf{W}) = (g^{-\top}) \mathbf{W},
\end{equation}
where $\mathbf{W}$ is a wrench (element of the dual algebra $\mathfrak{se}^*(3)$). In component form, if $\mathbf{W} = [\boldsymbol{\tau}; \mathbf{F}]$
(torque and force), then
\begin{equation}
  \text{Ad}^*_g \begin{bmatrix} \boldsymbol{\tau} \\ \mathbf{F} \end{bmatrix} = \begin{bmatrix} R^\top(\boldsymbol{\tau} - \mathbf{p} \times \mathbf{F}) \\ R^\top \mathbf{F} \end{bmatrix}.
  \label{eq:ch9:coadjoint-matrix}
\end{equation}

This transformation is essential when changing reference frames: if a wrench is applied at one point and we wish
to express it at another, the coadjoint action computes the equivalent wrench at the new frame.

\subsection{Poincaré Equations on SE(3)}

\index{Poincaré equations}%
For a rigid body with inertia tensor $I \in \mathbb{R}^{3 \times 3}$ and mass $m$, the kinetic energy
expressed in the body frame (where the frame is attached to the body) is
\begin{equation}
  T = \frac{1}{2}\xi^\top M_b \xi,
\end{equation}
where
\begin{equation}
  M_b = \begin{bmatrix} I & 0 \\ 0 & m I \end{bmatrix}
\end{equation}
is the body-frame mass matrix and $\xi$ is the body-frame twist.

The \textbf{Poincaré equations} describe the evolution of body-frame momentum:
\begin{equation}
  \frac{d}{dt}(M_b \xi) = \text{ad}^*_\xi(M_b \xi) + \mathbf{W}_{\text{ext}},
  \label{eq:ch9:poincare}
\end{equation}
where $\text{ad}^*_\xi$ is the coadjoint action of the infinitesimal motion $\xi$ and $\mathbf{W}_{\text{ext}}$
is the external wrench (force and torque).

In component form, the angular momentum $\boldsymbol{\ell} = I\boldsymbol{\omega}$ and linear momentum $\mathbf{p} = m\mathbf{v}$
evolve as
\begin{align}
  \dot{\boldsymbol{\ell}} &= \boldsymbol{\ell} \times \boldsymbol{\omega} + \boldsymbol{\tau}_{\text{ext}},\\
  \dot{\mathbf{p}} &= \mathbf{p} \times \boldsymbol{\omega} + \mathbf{F}_{\text{ext}}.
\end{align}

These are the generalized Euler equations for rigid body dynamics.

\subsection{Worked Example: Golf Grip Constraint in Twist Space}

\index{Golf swing!grip constraint}%
Consider a golfer gripping a club with both hands. Let the left hand contact the club at point $p_L$ and the right
hand at point $p_R$. The constraint is that both hands must hold the club without slipping, so their velocities
must equal the club's velocity at those points.

In twist space, the constraint Jacobian relates joint velocities to the twists of each contact point:
\begin{equation}
  J_{\text{left}} \dot{q}_{\text{left}} = \xi_{\text{club}} |_{p_L}, \quad J_{\text{right}} \dot{q}_{\text{right}} = \xi_{\text{club}} |_{p_R}.
\end{equation}

Eliminating $\xi_{\text{club}}$ (by requiring both equal the same club twist), we get a bilateral constraint:
\begin{equation}
  J_{\text{left}} \dot{q}_{\text{left}} = J_{\text{right}} \dot{q}_{\text{right}},
\end{equation}
or equivalently,
\begin{equation}
  [J_{\text{left}}, -J_{\text{right}}] \begin{bmatrix} \dot{q}_{\text{left}} \\ \dot{q}_{\text{right}} \end{bmatrix} = 0.
\end{equation}

The constraint Jacobian is $J_c = [J_{\text{left}}, -J_{\text{right}}] \in \mathbb{R}^{6 \times (n_L + n_R)}$.

The null space of $J_c$ parameterizes motions compatible with the grip: two arms moving in tandem while maintaining
contact. Any griping motion must lie in $\ker(J_c)$.

The constraint wrench space (reaction forces at the grip) is $\text{range}(J_c^\top)$. These are the forces and
torques the club exerts on the hands, which must be balanced by muscle forces in the arms.

\subsection{Grassmann Geometry and Singularities}

\index{Grassmann variety}%
\index{Singularity!Grassmann}%
The Grassmann variety $\text{Gr}(k, 6)$ is the space of $k$-dimensional linear subspaces of $\mathbb{R}^6$ (the twist space).
A set of $k$ joint screws (twist generators) span a $k$-dimensional subspace; the Grassmann variety parameterizes
all possible such subspaces.

For a parallel mechanism with $k$ joint screws $\xi_1, \ldots, \xi_k$, the Grassmann point is
\begin{equation}
  \Pi = \text{span}\{\xi_1, \ldots, \xi_k\} \in \text{Gr}(k, 6).
\end{equation}

A singularity occurs when this Grassmann point reaches a location where the screws become linearly dependent.
More precisely, at a singular configuration, the $k$ screws no longer span a $k$-dimensional subspace:
\begin{equation}
  \text{rank}([\xi_1, \ldots, \xi_k]) < k.
\end{equation}

The set of singular configurations forms a codimension-1 subvariety of the Grassmann variety (generically).

For the golf swing, singularities occur when the two arms align (both arms vertical, for example), making it impossible
to separately control the arm configuration while maintaining a grip. Near singularities, the Jacobian becomes
ill-conditioned, requiring large forces or torques to move the club.

\index{Loop closure}%
In a closed-loop kinematic chain (e.g., a parallel mechanism or a multi-loop robot), the constraint
is that the twists must satisfy a \textbf{loop closure} condition. If the chain starts at link 0,
passes through links $1, 2, \ldots, N$, and returns to link 0, then
\begin{equation}
  \sum_{i=1}^{N} \text{Ad}_{T_{i-1}} \xi_i = 0,
  \label{eq:ch9:loop-closure-twist}
\end{equation}
where $\xi_i$ is the twist of joint $i$ (relative motion between links $i-1$ and $i$),
$T_{i-1}$ transforms from link $i$ to the base, and $\text{Ad}$ is the adjoint action.

In differential form, this means the end-effector velocity is constrained by the joint velocities
of all limbs reaching it:
\begin{equation}
  J_1(q_1)\dot{q}_1 = J_2(q_2)\dot{q}_2 = \cdots = J_N(q_N)\dot{q}_N = \mathbf{v}_{\text{ee}},
  \label{eq:ch9:velocity-equality}
\end{equation}
where $J_i$ is the geometric Jacobian of limb $i$.

\subsection{Reciprocal Screws and Singular Configurations}

\index{Reciprocal screw}%
A screw $\xi$ is \textbf{reciprocal} to another screw $\eta$ if
\begin{equation}
  \xi^\top M_b \eta = 0.
\end{equation}

In a grasping problem, the \textbf{reciprocal wrench space} is the set of wrenches that produce zero
power on all achievable twists by the fingers. At a singularity, the reciprocal space becomes non-trivial,
meaning certain wrenches cannot be resisted.

\section{Optimal Control and Pontryagin Maximum Principle}
\label{sec:ch9:optimal-control}

\subsection{Optimal Control Problem Formulation}

Consider the constrained optimal control problem:
\begin{equation}
  \begin{aligned}
    \min_u \quad & \int_0^T \ell(q(t), \dot{q}(t), u(t)) \, dt + \phi(q(T)) \\
    \text{subject to} \quad & \dot{q} = q', \\
    & \ddot{q} = f(q,\dot{q}) + G(q)u, \\
    & \Phi(q) = 0, \\
    & u(t) \in \mathcal{U}.
  \end{aligned}
  \label{eq:ch9:optimal-control-problem}
\end{equation}

\index{Optimal control}%
where $\mathcal{U}$ is the set of admissible controls (typically $\mathcal{U} = [-u_{\max}, u_{\max}]^m$ or
$\mathcal{U} = \{u : \|u\| \le 1\}$).

\subsection{The Hamiltonian and Pontryagin Maximum Principle}

\index{Hamiltonian!Pontryagin}%
\index{Pontryagin maximum principle}%
Define the Hamiltonian:
\begin{equation}
  H(q, \dot{q}, \lambda_q, \lambda_{\dot{q}}, u) = \ell(q,\dot{q},u) + \lambda_q^\top \dot{q} + \lambda_{\dot{q}}^\top(f(q,\dot{q}) + G(q)u),
  \label{eq:ch9:hamiltonian}
\end{equation}
where $\lambda_q$ and $\lambda_{\dot{q}}$ are costates (Lagrange multipliers for the differential equations).

\begin{theorem}[Pontryagin Maximum Principle]\label{thm:pmp}
  If $(q^*, \dot{q}^*, u^*)$ is a solution to the optimal control problem \eqref{eq:ch9:optimal-control-problem},
  then there exist costates $\lambda_q(t)$ and $\lambda_{\dot{q}}(t)$ such that:
  \begin{enumerate}
    \item The costate equations hold:
    \begin{align}
      \dot{\lambda}_q &= -\frac{\partial H}{\partial q} + J_c(q)^\top \mu,\\
      \dot{\lambda}_{\dot{q}} &= -\frac{\partial H}{\partial \dot{q}},
    \end{align}
    where $\mu$ are Lagrange multipliers for the constraint $\Phi(q) = 0$.
    \item The optimal control minimizes the Hamiltonian:
    \begin{equation}
      u^* = \arg\min_{u \in \mathcal{U}} H(q^*, \dot{q}^*, \lambda_q, \lambda_{\dot{q}}, u).
    \end{equation}
    \item Boundary conditions are satisfied: $\lambda_q(T) = \nabla\phi(q(T))$.
  \end{enumerate}
\end{theorem}

\subsection{Switching Functions and Singular Arcs}

\index{Singular arc}%
\index{Switching function}%
When the control appears linearly in the Hamiltonian, i.e., $H = H_0 + (G(q)^\top \lambda_{\dot{q}})^\top u$,
the optimal control law is bang-bang (at the boundary of $\mathcal{U}$) unless the \textbf{switching function}
\begin{equation}
  \phi_i(t) := G_i(q(t))^\top \lambda_{\dot{q}}(t) = \frac{\partial H}{\partial u_i}
  \label{eq:ch9:switching-function}
\end{equation}
vanishes over an interval. When $\phi_i \ne 0$, the optimal control is
\begin{equation}
  u_i^* = \begin{cases} u_{\max} & \text{if } \phi_i(t) > 0, \\ u_{\min} & \text{if } \phi_i(t) < 0, \end{cases}
\end{equation}
corresponding to maximum (or minimum) control effort in that direction.

When $\phi_i(t) = 0$ over a time interval, the control is \textbf{singular}: it is determined not by minimization
but by maintaining the constraint $\phi_i = 0$. The singular control is found by requiring that $\dot{\phi}_i = 0$,
$\ddot{\phi}_i = 0$, etc., until $u$ appears explicitly.

For a golf swing, bang-bang control would correspond to muscles either fully contracted or fully relaxed at each instant.
Biological evidence suggests that actual swings use \emph{smooth} control profiles, with muscle activation varying
continuously. This indicates either (1) a different cost function (not time-optimal), or (2) constraints on the
rate of muscle force change (not instantaneous activation).

\subsection{Quadratic Cost and Optimal Control}

For a quadratic cost functional
\begin{equation}
  J = \int_0^T \left( \|q(t) - q_d\|_Q^2 + \|u(t)\|_R^2 \right) dt,
\end{equation}
where $Q$ is a state weight and $R$ is a control weight, the Hamiltonian becomes
\begin{equation}
  H = \|q - q_d\|_Q^2 + \|u\|_R^2 + \lambda_q^\top \dot{q} + \lambda_{\dot{q}}^\top(f(q,\dot{q}) + G(q)u).
\end{equation}

The optimal control minimizes $H$ with respect to $u$:
\begin{equation}
  \frac{\partial H}{\partial u} = 2R u + G(q)^\top \lambda_{\dot{q}} = 0,
\end{equation}
giving
\begin{equation}
  u^* = -\frac{1}{2} R^{-1} G(q)^\top \lambda_{\dot{q}}.
  \label{eq:ch9:optimal-control-quadratic}
\end{equation}

This is a linear feedback law in terms of the costate. For a golf swing with bounded actuation (e.g., $u_i \in [0, 1]$
for muscle activation), the control clips at the boundaries:
\begin{equation}
  u_i^* = \text{clip}\left[ -\frac{1}{2} (R^{-1} G^\top \lambda_{\dot{q}})_i, 0, 1 \right],
  \label{eq:ch9:clipped-control}
\end{equation}
a mixed bang-singular control.

\subsection{Hamilton-Jacobi-Bellman Equation}

\index{Hamilton-Jacobi-Bellman}%
The optimal value function $V(q, \dot{q}, t)$ (minimum cost from state $(q, \dot{q})$ at time $t$) satisfies
the HJB equation:
\begin{equation}
  -\frac{\partial V}{\partial t} = \min_u \left[ \ell(q,\dot{q},u) + \frac{\partial V}{\partial q}^\top \dot{q}
  + \frac{\partial V}{\partial \dot{q}}^\top(f + Gu) \right],
  \label{eq:ch9:hjb}
\end{equation}
with boundary condition $V(q, \dot{q}, T) = \phi(q)$.

The HJB equation is a nonlinear first-order PDE. Its solution encodes the optimal control law everywhere in state space:
\begin{equation}
  u^*(q, \dot{q}, t) = \arg\min_u \left[ \ell(q,\dot{q},u) + \frac{\partial V}{\partial \dot{q}}^\top G(q)u \right].
\end{equation}

For systems without explicit time dependence, the steady-state HJB is
\begin{equation}
  0 = \min_u \left[ \ell(q,\dot{q},u) + \frac{\partial V}{\partial q}^\top \dot{q}
  + \frac{\partial V}{\partial \dot{q}}^\top(f + Gu) \right].
\end{equation}

\subsubsection{Curse of Dimensionality}

The HJB equation must be solved over the full state space $(q, \dot{q})$. For a human golf swing:
\begin{itemize}
  \item Configuration space: $q$ has dimension $n \approx 14$ (7 joint angles per arm + torso).
  \item Tangent bundle: $(q, \dot{q})$ has dimension $2n \approx 28$.
  \item Discretization: If we discretize each dimension with 10 grid points, we have $10^{28}$ grid cells---computationally intractable.
\end{itemize}

This is the \textbf{curse of dimensionality}: the number of grid points grows exponentially with state dimension.
Direct numerical solution of HJB is thus infeasible for high-dimensional systems.

\subsubsection{Approximate Methods}

To overcome this limitation:
\begin{enumerate}
  \item \textbf{iLQR/DDP} (Iterative Linear-Quadratic Regulator / Differential Dynamic Programming): Linearize
  dynamics around a nominal trajectory and solve a sequence of finite-horizon, quadratic approximations.
  This is tractable even for 28-dimensional systems.

  \item \textbf{Value Iteration with Basis Functions}: Approximate $V(q, \dot{q}) \approx \sum_k w_k \phi_k(q, \dot{q})$
  using learned basis functions (e.g., radial basis functions or neural networks). The HJB equation becomes
  a regression problem.

  \item \textbf{Policy Search}: Parameterize the control law $u = \pi(q, \dot{q}; \theta)$ with parameters $\theta$,
  and optimize directly over $\theta$ using gradient descent or evolutionary strategies.
\end{enumerate}

For biomechanics, evidence suggests that humans do not solve the global optimization problem online. Rather,
the nervous system likely uses learned inverse models (neural networks) or heuristic strategies (e.g., minimizing
energy or jerk) that approximate the optimal solution.

\subsection{iLQR and Trajectory Optimization}

\index{iLQR}%
\index{DDP}%
Given a nominal trajectory $\{(q_k, \dot{q}_k, u_k)\}_{k=0}^T$, the iLQR (iterative linear-quadratic regulator)
algorithm:
\begin{enumerate}
  \item \textbf{Backward pass}: Linearize the dynamics around the nominal trajectory:
  \begin{equation}
    \delta\ddot{q}_k = A_k \delta q_k + B_k \delta u_k,
  \end{equation}
  Quadratize the cost around the trajectory, and solve the Riccati recursion backward to compute
  optimal feedback gains $K_k$ and value-function matrices $S_k$.
  \item \textbf{Forward pass}: Apply the feedback law $\delta u_k = K_k \delta q_k$ and integrate forward
  to compute a new nominal trajectory.
  \item \textbf{Repeat} until the trajectory converges.
\end{enumerate}

When the constraints must be satisfied, we augment the backward pass by projecting both the linearized
dynamics and the cost gradient onto the constraint manifold.

\section{Observability, Identifiability, and Inverse Problems}
\label{sec:ch9:observability}

\subsection{Observability Rank Condition}

\index{Observability}%
\index{Observability rank matrix}%
A system is \textbf{observable from state} $(q, \dot{q})$ with measurements $y = h(q, \dot{q})$ if
the mapping from initial states to output trajectories is injective. The observable subspace
can be characterized via the \textbf{observability rank condition} (Hermann \& Krener 1977).

For a control-affine system
\begin{equation}
  \dot{\mathbf{x}} = f(\mathbf{x}) + G(\mathbf{x})u, \quad y = h(\mathbf{x}),
\end{equation}
define the observability rank matrix
\begin{equation}
  \mathcal{O}(\mathbf{x}) = \begin{bmatrix}
    \frac{\partial h}{\partial \mathbf{x}} \\
    \frac{\partial (L_f h)}{\partial \mathbf{x}} \\
    \frac{\partial (L_f^2 h)}{\partial \mathbf{x}} \\
    \vdots
  \end{bmatrix},
  \label{eq:ch9:observability-matrix}
\end{equation}
where $L_f h$ is the Lie derivative (directional derivative along $f$): $L_f h = \frac{\partial h}{\partial \mathbf{x}} f(\mathbf{x})$.

Equivalently, define the observability codistribution:
\begin{equation}
  d\mathcal{O} = \text{span}\{dh, L_f dh, L_f^2 dh, \ldots, L_f^{n-1} dh, L_{g_i} dh, L_{g_i} L_f dh, \ldots\},
\end{equation}
where $L_X dh$ is the Lie derivative of the differential form $dh$ along vector field $X$.

\begin{proposition}[Observability from Measurements]\label{prop:observability}
  If the observability rank matrix $\mathcal{O}(\mathbf{x})$ has full rank (rank $n$) at all points
  or on an open dense set, then the system is (generically) locally observable.
\end{proposition}

\subsection{Drift Force Observability in Mechanical Systems}

\index{Drift force}%
\begin{proposition}[Drift Force Observability in Constrained Mechanical Systems]\label{prop:drift-observability}
  Consider a mechanical system with constrained dynamics
  \begin{equation}
    M(q)\ddot{q} + C(q,\dot{q})\dot{q} + \nabla V(q) = J_c(q)^\top\lambda + B(q)u,
    \quad \Phi(q) = 0.
  \end{equation}
  If measurements consist of $q(t)$ and $\dot{q}(t)$ (position and velocity), then the drift field
  \begin{equation}
    f(q,\dot{q}) = -M(q)^{-1}(C(q,\dot{q})\dot{q} + \nabla V(q) - J_c(q)^\top\lambda)
  \end{equation}
  is \textbf{fully observable}, because $M$, $C$, $\nabla V$, and $J_c$ all depend only on measurable
  quantities $(q, \dot{q})$.
\end{proposition}

\begin{proof}
  By assumption, $M(q)$, $C(q,\dot{q})$, and $\nabla V(q)$ are known functions of $(q, \dot{q})$.
  The constraint forces $J_c(q)^\top\lambda$ are also determined by the constraint equations and
  the observed $(q, \dot{q})$: the acceleration constraint $J_c(q)\ddot{q} + \dot{J}_c(q,\dot{q})\dot{q} = 0$
  determines $\ddot{q}$ uniquely (given $q, \dot{q}$), and hence $\lambda$ follows from
  \begin{equation}
    \lambda = (J_c(q)M(q)^{-1}J_c(q)^\top)^{-1} J_c(q)M(q)^{-1}(B(q)u - C(q,\dot{q})\dot{q} - \nabla V(q)).
  \end{equation}
  Therefore, $f(q,\dot{q})$ is fully determined from measurements.
\end{proof}

\subsection{Identifiability Ambiguity and the Observable Subspace}

\index{Identifiability}%
\index{Ambiguity set}%
In practice, we may not measure all components of $(q,\dot{q})$. If measurements are $y = h(q, \dot{q})$,
then only the observable subspace is identifiable:
\begin{equation}
  \mathcal{O}_{\text{obs}} = \ker(d\mathcal{O})^{\perp}.
\end{equation}

States in the unobservable subspace $\ker(d\mathcal{O})$ produce identical measurement trajectories.
The \textbf{ambiguity set} at state $\mathbf{x}_0$ is an affine subspace:
\begin{equation}
  \mathcal{A}(\mathbf{x}_0) = \{\mathbf{x}_0 + \mathbf{v} : \mathbf{v} \in \ker(d\mathcal{O}(\mathbf{x}_0))\},
\end{equation}
a coset of the null space. All points in this set are indistinguishable from measurements alone.

For biomechanical systems in gait analysis, the measurements typically include:
\begin{itemize}
  \item Joint angles $(q)$ from motion capture.
  \item Joint angular velocities $(\dot{q})$ from pose derivatives.
  \item Ground reaction forces (from force plates).
  \item (Optionally) EMG from muscles.
\end{itemize}

The unobservable subspace typically includes internal muscle activations and muscle-tendon parameters (force-length
and force-velocity relationships). Different muscle activation patterns can produce the same motion and GRF,
creating ambiguity.

This identifiability problem is resolved via:
\begin{enumerate}
  \item \textbf{Electromyography (EMG)}: Direct measurement of muscle electrical activity. EMG correlates with muscle force
  (though nonlinearly, with a temporal delay ~50 ms), providing additional observational data. Integration over
  time gives muscle activation (a control input), which constrains the null-space ambiguity.

  \item \textbf{Optimization criteria}: In the absence of EMG, assume the nervous system optimizes a physiological criterion.
  Common choices:
  \begin{itemize}
    \item \textbf{Minimum effort}: Minimize $\sum_i u_i^2$ (control effort).
    \item \textbf{Minimum metabolic cost}: Minimize $\sum_i C_i(u_i)$ where $C_i$ is the metabolic cost of muscle $i$
    \cite{Zajac1989}.
    \item \textbf{Minimum fatigue}: Minimize $\sum_i u_i^p / A_i$ (muscle stress), as in the Crowninshield \& Brand criterion
    \cite{CrowninshieldBrand1981}.
  \end{itemize}

  \item \textbf{Instrumented equipment}: Pressure sensors in shoes, inertial measurement units (IMUs), and wireless
  EMG systems provide additional measurements that constrain the unobservable subspace.
\end{enumerate}

The choice of optimization criterion strongly influences the predicted muscle activations. Biomechanical analyses
often compare predictions from different criteria to test physiological plausibility \cite{TodorovJordan2002}.

\section{Ground Reaction Forces and Lower-Body Parallel Mechanisms}
\label{sec:ch9:grf}

\subsection{The Bilateral Leg Mechanism}

\index{Ground reaction force}%
\index{GRF}%
\index{Parallel mechanism!lower body}%
The human (or biped) lower body can be viewed as a parallel mechanism: two legs, each a serial chain,
meeting at the pelvis and terminating in feet on the ground. When both feet are in contact, the system
becomes constrained.

The unconstrained DOF of each leg is 7 (3 DOF hip, 1 DOF knee, 3 DOF ankle). For both legs:
\begin{equation}
  \text{Total DOF}_{\text{free}} = 2 \times 7 = 14.
\end{equation}

When one foot is in contact with the ground (single-leg stance), the foot position and orientation are fixed relative
to the ground, imposing 6 constraints (3 translational + 3 rotational). Thus:
\begin{equation}
  \text{DOF}_{\text{single-leg}} = 14 - 6 = 8,
\end{equation}
distributed as: 6 DOF for the pelvis/torso and 2 DOF for the free leg (internal configuration).

When both feet are in bilateral ground contact, we have additional constraints:
\begin{itemize}
  \item \textbf{Foot-ground contact}: Each foot contact point has zero velocity (no penetration, no slip):
  $2 \times 3 = 6$ translational constraints.
  \item \textbf{Foot orientation}: If we assume feet remain flat and parallel to the ground, each foot has
  3 constraints on rotation: $2 \times 3 = 6$ rotational constraints.
\end{itemize}

However, when both feet are on the ground, the feet may not be independently aligned (e.g., they share the ground plane).
The actual number of independent constraints for bilateral contact is typically:
\begin{equation}
  \text{Constraints}_{\text{bilateral}} \approx 6 + 6 = 12.
\end{equation}

Thus:
\begin{equation}
  \text{DOF}_{\text{bilateral}} = 14 - 12 = 2,
\end{equation}
leaving approximately 2 DOF for the system (e.g., height of the pelvis and one internal leg angle), though this
depends on foot geometry and friction assumptions.

\subsection{Force Plate Measurements and Wrench Decomposition}

\index{Force plate}%
A force plate mounted under the ground measures the reaction wrench at the contact surface. The six components are:
\begin{align}
  F_x, F_y &: \text{Horizontal forces (anterior-posterior and medial-lateral)}, \\
  F_z &: \text{Vertical normal force}, \\
  M_x, M_y, M_z &: \text{Moments (torques)} \text{ about the x, y, z axes}.
\end{align}

Additionally, the force plate provides:
\begin{equation}
  \text{COP} = (\text{COP}_x, \text{COP}_y) := \left( \frac{M_y}{F_z}, -\frac{M_x}{F_z} \right),
\end{equation}
the center of pressure (the location on the ground where a single equivalent force acts). This is computed from
the moments assuming all forces act through a single point on the ground surface.

\subsection{Sequential Inverse Dynamics Computation}

\index{Inverse dynamics!sequential}%
Given measured kinematics $(q(t), \dot{q}(t), \ddot{q}(t))$ and GRF $\mathbf{W}_{\text{grf}}$, the inverse dynamics
computes joint torques using a backward recursion from the foot to the hip.

\subsubsection{Newton-Euler Recursion}

For each link $i$ (starting from the end-effector at the foot):
\begin{enumerate}
  \item Compute the acceleration of the link's center of mass:
  \begin{equation}
    \mathbf{a}_{i,\text{cm}} = \mathbf{a}_{i} + \boldsymbol{\alpha}_{i} \times \mathbf{r}_{i,\text{cm}} - \boldsymbol{\omega}_{i} \times (\boldsymbol{\omega}_{i} \times \mathbf{r}_{i,\text{cm}}),
  \end{equation}
  where $\mathbf{a}_i$ and $\boldsymbol{\alpha}_i$ are the linear and angular acceleration of link $i$, $\mathbf{r}_{i,\text{cm}}$
  is the position of the center of mass relative to the joint, and $\boldsymbol{\omega}_i$ is angular velocity.

  \item Compute the force and moment at link $i$:
  \begin{align}
    \mathbf{f}_i &= \mathbf{f}_{i+1} + m_i \mathbf{a}_{i,\text{cm}}, \label{eq:ch9:force-recursion} \\
    \boldsymbol{\tau}_i &= \boldsymbol{\tau}_{i+1} + \mathbf{r}_{i,\text{cm}} \times \mathbf{f}_i + I_i \boldsymbol{\alpha}_i + \boldsymbol{\omega}_i \times (I_i \boldsymbol{\omega}_i),
  \end{align}
  where $m_i$ is the link mass, $I_i$ is the inertia tensor, and the last term is the gyroscopic effect.

  \item Extract the joint torque as the component along the joint axis:
  \begin{equation}
    \tau_i = \boldsymbol{\tau}_i \cdot \hat{e}_i,
  \end{equation}
  where $\hat{e}_i$ is the unit vector along the joint axis.
\end{enumerate}

The boundary condition at the foot is the measured GRF: $\mathbf{f}_{\text{foot}} = \mathbf{F}_{\text{grf}}$ and
$\boldsymbol{\tau}_{\text{foot}} = \mathbf{M}_{\text{grf}}$ (or computed from the COP).

This backward recursion is efficient ($O(n)$ complexity) and does not require matrix inversion, in contrast to
forward dynamics.

\subsection{Kinetic Energy Partition and Power Balance}

\index{Power balance}%
\index{Mechanical work}%
The total kinetic energy of the lower body is partitioned among links:
\begin{equation}
  T_{\text{total}} = T_{\text{pelvis}} + T_{\text{left-leg}} + T_{\text{right-leg}} + \ldots
\end{equation}

Each link contributes:
\begin{equation}
  T_i = \frac{1}{2} m_i v_{i,\text{cm}}^2 + \frac{1}{2} \boldsymbol{\omega}_i^\top I_i \boldsymbol{\omega}_i.
\end{equation}

The mechanical power delivered by the ground at each contact point is:
\begin{equation}
  P_{\text{grf}} = \mathbf{F}_{\text{grf}} \cdot \mathbf{v}_{\text{contact}} + \mathbf{M}_{\text{grf}} \cdot \boldsymbol{\omega}_{\text{foot}}.
\end{equation}

For a stationary foot (no slipping, $\mathbf{v}_{\text{contact}} = 0$), the power is purely rotational:
\begin{equation}
  P_{\text{grf}} = \mathbf{M}_{\text{grf}} \cdot \boldsymbol{\omega}_{\text{foot}}.
  \label{eq:ch9:ground-torque-power}
\end{equation}

This ground torque power is often overlooked in simplified models but is essential for accurate energy accounting:
it represents rotational energy supplied by the ground, critical for activities like cutting or rapid direction changes.

The power balance (without friction losses) is:
\begin{equation}
  P_{\text{grf}} + P_{\text{muscles}} = \frac{d}{dt}(T_{\text{total}} + V_{\text{total}}),
\end{equation}
where $P_{\text{muscles}} = \sum_j \tau_j \dot{q}_j$ is the mechanical power generated by muscles at each joint.

\subsection{Ground Wrench Decomposition}

\index{Wrench decomposition!GRF}%
The wrench exerted by the ground on a foot is $\mathbf{W} = (\mathbf{F}, \boldsymbol{\tau})$, consisting of:
\begin{itemize}
  \item \textbf{Normal force} $F_z$: Perpendicular to the ground, acts to prevent penetration.
  \item \textbf{Horizontal friction forces} $F_x, F_y$: In the plane of the ground.
  \item \textbf{Moment} $\boldsymbol{\tau} = (\tau_x, \tau_y, \tau_z)$: The torque transmitted through the contact.
\end{itemize}

For modeling, we often decompose the wrench into components aligned with the contact surface. If the
contact region has a pressure distribution $p(\xi, \eta)$ over area $A$, the total wrench is
\begin{align}
  F_z &= \int_A p(\xi,\eta) \, d\xi d\eta, \\
  \mathbf{F}_{xy} &= \text{(friction traction from constitutive model)}, \\
  \boldsymbol{\tau} &= \int_A (\xi, \eta, 0) \times \begin{bmatrix} F_x(\xi,\eta) \\ F_y(\xi,\eta) \\ p(\xi,\eta) \end{bmatrix} d\xi d\eta.
\end{align}

\subsection{Inverse Dynamics with GRF as Boundary Condition}

\index{Inverse dynamics}%
Given a measured motion $(q(t), \dot{q}(t), \ddot{q}(t))$ and a force plate measurement $\mathbf{W}_{\text{grf}}(t)$,
the inverse dynamics problem seeks the joint torques $\tau(t)$.

For a serial chain (single leg) with a base trajectory, the Newton-Euler recursion proceeds:
\begin{enumerate}
  \item Initialize the end-effector acceleration using the measured acceleration and the GRF.
  \item Propagate backward through the kinematic chain using the recursive formula:
  \begin{equation}
    \mathbf{f}_i = \mathbf{f}_{i+1} + m_i \mathbf{a}_{i,\text{cm}}, \quad \boldsymbol{\tau}_i = \boldsymbol{\tau}_{i+1} + \mathbf{r}_{i,\text{cm}} \times \mathbf{f}_i + I_i \boldsymbol{\alpha}_i,
  \end{equation}
  where $\mathbf{f}_i$ and $\boldsymbol{\tau}_i$ are force and torque transmitted through joint $i$,
  $\mathbf{a}_{i,\text{cm}}$ and $\boldsymbol{\alpha}_i$ are acceleration and angular acceleration of the link center of mass,
  $m_i$ and $I_i$ are mass and inertia tensor.
  \item The joint torque is $\tau_i = \boldsymbol{\tau}_i \cdot \hat{e}_i$, the component along the joint axis.
\end{enumerate}

This backward recursion is efficient and does not require matrix inversion (unlike forward dynamics).

\subsection{Kinetic Energy Partition and Ground Torque}

\index{Ground torque}%
The power delivered by the ground wrench is
\begin{equation}
  P_{\text{grf}} = \mathbf{F}^\top \mathbf{v}_{\text{contact}} + \boldsymbol{\tau}^\top \boldsymbol{\omega}_{\text{foot}}.
\end{equation}

For a stationary foot (no slipping), $\mathbf{v}_{\text{contact}} = 0$, so all mechanical power input from the
ground is rotational:
\begin{equation}
  P_{\text{grf}} = \boldsymbol{\tau}^\top \boldsymbol{\omega}_{\text{foot}}.
\end{equation}

This torque is critical: it accounts for the rotational contribution of the ground to the system's kinetic
energy, which is often overlooked in simplified biomechanical models.

\section{Connections, Curvature, and Geometric Phase}
\label{sec:ch9:connections}

\subsection{Mechanical Connections on Principal Bundles}

\index{Mechanical connection}%
When a system has symmetry (e.g., invariance under a group action $G$), we can decompose the configuration
space as a principal bundle $\pi: Q \to Q/G$ with structure group $G$.

A \textbf{mechanical connection} $A: TQ \to \mathfrak{g}$ (the Lie algebra of $G$) decomposes velocities:
\begin{equation}
  \dot{q} = \dot{q}_{\text{shape}} + \text{(group velocity)},
\end{equation}
where $\dot{q}_{\text{shape}}$ is the rate of change of shape (on the quotient $Q/G$) and the group velocity
is the rate of change of the group variable.

For a system with a cyclic symmetry (e.g., a falling cat that can change its orientation without external
torque), the connection is provided by the constraint that angular momentum is conserved in the absence
of external torques.

\subsection{Locked Inertia Tensor and Momentum Map}

\index{Momentum map}%
\index{Inertia tensor!locked}%
\index{Noether's theorem}%
If we ``lock'' the shape variables (i.e., hold $q_{\text{shape}}$ fixed), the kinetic energy in the group
direction defines the \textbf{locked inertia tensor} $I_{\text{locked}}(q_{\text{shape}})$.

The \textbf{momentum map} is the function $\mathbf{J}: T Q \to \mathfrak{g}^*$ (dual of the Lie algebra) that
extracts the generalized momentum in the group direction. For a system with a Lagrangian that is invariant
under $G$-action, Noether's theorem implies that the momentum map is conserved along trajectories:
\begin{equation}
  \mathbf{J}(q(t), \dot{q}(t)) = \text{constant}.
\end{equation}

\subsection{Geometric Phase and Holonomy}

\index{Geometric phase}%
\index{Holonomy}%
\index{Berry phase}%
Consider a system that undergoes a cyclic shape change: $q_{\text{shape}}(0) = q_{\text{shape}}(T)$ but
the group variable changes by $g(T) \ne \text{id}$. The change in group variable is the \textbf{geometric phase}
(also called Berry phase or Hannay angle, after Berry (1984) and Hannay (1985); see also Shapere \& Wilczek (1989)).

\begin{example}[Geometric Phase in a Golf Swing]
  A golfer swinging both arms (shape change) in a coordinated pattern rotates the torso relative to the pelvis
  (group change). Even without external torques (conserving total angular momentum), the asymmetry and sequencing
  of arm motion produces a net rotation. This is a geometric phase effect: certain shape changes generate rotation.

  Specifically, if the golfer:
  \begin{enumerate}
    \item Raises the left arm (rotates about the shoulder).
    \item Lowers the left arm while raising the right arm.
    \item Lowers the right arm.
    \item (Returns to initial shape, but torso has rotated.)
  \end{enumerate}

  The sequence matters: a different order of arm raises produces a different net torso rotation. This is purely
  geometric (due to the Lie group structure), not from any external torque.

  The rate of accumulation of geometric phase is proportional to the area enclosed by the shape trajectory
  in the shape space, weighted by the inverse locked inertia tensor: $\Delta\phi \propto \oint A(q_{\text{shape}}) \cdot dq_{\text{shape}}$.
\end{example}

\subsubsection{Mechanical Connection and Berry Phase Formula}

The \textbf{mechanical connection} is the projection $A(q, \dot{q}): TQ \to \mathfrak{g}$ that decomposes velocities into
shape and group parts. It is defined via the locked inertia:
\begin{equation}
  A(q_{\text{shape}}, \dot{q}_{\text{shape}}) := I_{\text{locked}}(q_{\text{shape}})^{-1} J_G(q_{\text{shape}}, \dot{q}_{\text{shape}}),
  \label{eq:ch9:mechanical-connection}
\end{equation}
where $J_G$ is the momentum map (extracts the group-direction component of momentum from shape changes).

The accumulated geometric phase along a closed curve in shape space is given by:
\begin{equation}
  \Delta \phi = \oint A(q_{\text{shape}}) \cdot dq_{\text{shape}} = \int_{\text{interior}} F(q_{\text{shape}}) \, dA,
  \label{eq:ch9:berry-phase-formula}
\end{equation}
where $F$ is the curvature 2-form (the exterior derivative of $A$), and the second integral is over the surface
bounded by the closed curve.

\subsubsection{Grassmann Geometry and Manipulability}

\index{Manipulability!measure}%
The \textbf{manipulability measure} of a Jacobian $J(q)$ quantifies how well the system can apply forces in all
directions:
\begin{equation}
  w(q) = \sqrt{\det(J(q) J(q)^\top)}.
  \label{eq:ch9:manipulability}
\end{equation}

High manipulability indicates that the Jacobian has full rank and is well-conditioned; low manipulability indicates
near-singularity.

The connection to curvature: manipulability is related to the volume element in the Grassmann variety. When the
joint screws are nearly parallel (low manipulability), the curvature of the mechanical connection becomes large,
meaning the geometric phase changes rapidly with shape variations. This suggests that near singularities, small
shape changes produce large rotational effects---a useful property for certain tasks (like rapid reorientation)
but problematic for precise control.

\section{Numerical Methods and Implementation}
\label{sec:ch9:numerics}

\subsection{SHAKE and RATTLE Methods}

\index{SHAKE}%
\index{RATTLE}%
\index{Constraint-preserving integrator}%
For systems with holonomic constraints, the most commonly used constraint-preserving integrators are:

\subsubsection{SHAKE Algorithm}

\textbf{SHAKE} (Ryckaert et al. 1977) is a position-level constraint solver that uses iterative Lagrange
multiplier updates. The algorithm:
\begin{enumerate}
  \item Predict position: $q_{n+1}^{\text{pred}} = q_n + \Delta t \, \dot{q}_n + \frac{\Delta t^2}{2} M(q_n)^{-1} F(q_n, \dot{q}_n)$.
  \item Correct for constraints iteratively:
  \begin{equation}
    q_{n+1}^{(k+1)} = q_{n+1}^{(k)} - J_c(q_{n+1}^{(k)})^\top (J_c(q_{n+1}^{(k)}) J_c(q_{n+1}^{(k)})^\top)^{-1} \Phi(q_{n+1}^{(k)}),
    \label{eq:ch9:shake-update}
  \end{equation}
  repeating until $\|\Phi(q_{n+1}^{(k)})\| < \epsilon_{\text{tol}}$ (constraint violation is small).
  \item Compute velocities: $\dot{q}_{n+1} = \frac{q_{n+1} - q_n}{\Delta t}$.
\end{enumerate}

Advantages: Simple to implement, preserves constraint manifold up to numerical precision.
Disadvantages: Velocity constraint ($J_c \dot{q} = 0$) is not explicitly enforced, leading to potential velocity drift.

\subsubsection{RATTLE Algorithm}

\textbf{RATTLE} (Andersen 1983) augments SHAKE by also enforcing the velocity constraint. The algorithm:
\begin{enumerate}
  \item Predict position: $q_{n+1}^{\text{pred}} = q_n + \Delta t \, \dot{q}_n + \frac{\Delta t^2}{2} M(q_n)^{-1} F_n$.
  \item Correct position (same as SHAKE) until $\Phi(q_{n+1}) \approx 0$.
  \item Predict velocity: $\dot{q}_{n+1}^{\text{pred}} = \dot{q}_n + \Delta t \, M(q_n)^{-1} (F_n + J_c(q_{n+1})^\top \lambda_n)$.
  \item Correct velocity:
  \begin{equation}
    \dot{q}_{n+1} = \dot{q}_{n+1}^{\text{pred}} - J_c(q_{n+1})^\top (J_c(q_{n+1}) M(q_{n+1})^{-1} J_c(q_{n+1})^\top)^{-1}
    J_c(q_{n+1}) \dot{q}_{n+1}^{\text{pred}},
  \end{equation}
  ensuring $J_c(q_{n+1}) \dot{q}_{n+1} = 0$.
\end{enumerate}

Advantages: Preserves both position and velocity constraints; superior energy conservation.
Disadvantages: More expensive (requires solving two linear systems per step).

Both methods require solving linear systems with the constraint Jacobian, which can be accelerated using:
\begin{itemize}
  \item \textbf{Direct solvers} (Cholesky, LU): Fast for dense or well-structured matrices.
  \item \textbf{Iterative solvers} (Conjugate Gradient, GMRES, Krylov methods): Efficient for sparse matrices.
\end{itemize}

\subsection{SVD for Null-Space Projection}

\index{SVD!null space}%
To project accelerations onto the constraint manifold, compute the singular value decomposition (SVD):
\begin{equation}
  J_c(q) = U \Sigma V^\top,
  \label{eq:ch9:svd-jc}
\end{equation}
where $U \in \mathbb{R}^{c \times c}$ and $V \in \mathbb{R}^{n \times n}$ are orthogonal matrices, and $\Sigma \in \mathbb{R}^{c \times n}$
is diagonal with singular values $\sigma_1 \ge \sigma_2 \ge \cdots \ge \sigma_c > 0$.

The columns of $V$ can be partitioned as $V = [V_{\text{range}}, V_{\text{null}}]$:
\begin{itemize}
  \item $V_{\text{range}} = V_{:,1:c}$: The first $c$ columns (corresponding to non-zero singular values).
  These span the range of $J_c^\top$, i.e., directions that increase constraint violation.
  \item $V_{\text{null}} = V_{:,c+1:n}$: The last $n - c$ columns (corresponding to zero singular values).
  These span the null space $\ker(J_c)$, i.e., directions that preserve constraints.
\end{itemize}

The null-space projector is:
\begin{equation}
  P_{\ker(J_c)} = V_{\text{null}} V_{\text{null}}^\top = I - V_{\text{range}} V_{\text{range}}^\top.
  \label{eq:ch9:null-projector}
\end{equation}

To project an acceleration onto the constraint manifold:
\begin{equation}
  \ddot{q}_{\text{proj}} = P_{\ker(J_c)} \ddot{q} = (I - V_{\text{range}} V_{\text{range}}^\top) \ddot{q},
\end{equation}
ensuring $J_c(q) \ddot{q}_{\text{proj}} = 0$ (acceleration constraint is satisfied).

The SVD is numerically stable and does not require matrix inversion. For large-scale systems, sparse SVD
algorithms (using Krylov subspace methods) can accelerate computation. The SVD also naturally reveals the
rank of $J_c$ via the number of large singular values; if $\sigma_c$ is small, $J_c$ is nearly rank-deficient,
indicating proximity to a singularity.

\subsection{Berry Phase Computation}

\index{Berry phase!computation}%
\index{Line integral!connection form}%
To compute the net geometric phase accumulated over a shape trajectory $q_{\text{shape}}(s)$ for $s \in [0,1]$,
we use the formula (Shapere \& Wilczek 1989):
\begin{equation}
  \theta_{\text{geometric}} = \oint \mathbf{A}(q_{\text{shape}}) \cdot dq_{\text{shape}},
  \label{eq:ch9:berry-phase-line-integral}
\end{equation}
where $\mathbf{A}$ is the connection form. By Stokes' theorem, this equals
\begin{equation}
  \theta_{\text{geometric}} = \int_{\text{Surface}} (\text{curl} \, \mathbf{A}) \, dA = \int_{\text{Surface}} \Omega \, dA,
\end{equation}
where $\Omega = d\mathbf{A}$ is the curvature 2-form.

The mechanical connection relates to the locked inertia:
\begin{equation}
  \mathbf{A}(q_{\text{shape}}) = I_{\text{locked}}(q_{\text{shape}})^{-1} \frac{\partial I_{\text{locked}}}{\partial q_{\text{shape}}}.
\end{equation}

Numerical implementation:
\begin{enumerate}
  \item \textbf{Grid construction}: Discretize the shape space on a fine grid. For a 2-DOF shape space, this is
  an $N \times N$ grid of shape values.

  \item \textbf{Locked inertia computation}: For each grid point $q_{\text{shape}}$:
  \begin{equation}
    I_{\text{locked}}(q_{\text{shape}}) = \frac{\partial T}{\partial \dot{q}_{\text{group}}^2}\bigg|_{q_{\text{shape}}, \dot{q}_{\text{group}}}
    \approx \frac{1}{2} \frac{\partial^2 (T - V)}{\partial \dot{q}_{\text{group}}^2},
  \end{equation}
  computed by evaluating the kinetic energy in the group direction while holding shape fixed.

  \item \textbf{Connection form}: Compute the connection form via finite differences:
  \begin{equation}
    A_i(q_{\text{shape}}) \approx \frac{I_{\text{locked}}(q_{\text{shape}} + \delta_i) - I_{\text{locked}}(q_{\text{shape}} - \delta_i)}{2 \delta q_{\text{shape}, i}}.
  \end{equation}

  \item \textbf{Line integral}: Integrate around a closed curve in shape space using the trapezoidal rule:
  \begin{equation}
    \oint \mathbf{A} \cdot dq_{\text{shape}} \approx \sum_{k} \mathbf{A}(q_{\text{shape}, k}) \cdot (q_{\text{shape}, k+1} - q_{\text{shape}, k}).
  \end{equation}

  \item \textbf{Stokes integral (alternative)}: Compute the curvature 2-form $\Omega = d\mathbf{A}$ and integrate
  over the interior surface using Simpson's rule or other quadrature.
\end{enumerate}

For a golf swing, this allows prediction of the net torso rotation given a sequence of arm raises and lowers.

\section{Integration: From Theory to Practice}
\label{sec:ch9:integration}

\subsection{Conceptual Mapping of Key Objects}

Table~\ref{tab:ch9:mapping} presents a unified view of how key concepts in differential geometry relate to
mechanical structure.

\begin{table}[h]
\centering
\small
\begin{tabular}{|p{2.5cm}|p{2cm}|p{2.5cm}|p{3cm}|}
\hline
\textbf{Differential Geometry} & \textbf{Mechanics} & \textbf{Example} & \textbf{Role}\\
\hline
Configuration manifold $Q$ & Rigid body positions & $SE(3)$ for 6-DOF body & State space \\
\hline
Riemannian metric $g = M$ & Kinetic energy & $T = \frac{1}{2}\dot{q}^\top M \dot{q}$ & Energy cost of motion\\
\hline
Christoffel symbols $\Gamma$ & Coriolis/centrifugal & $C(q,\dot{q})\dot{q}$ & Curved-space inertia \\
\hline
Constraint submanifold $Q_c$ & Contact/joint limits & $\Phi(q) = 0$ & Feasible configuration set \\
\hline
Constraint Jacobian $J_c$ & Constraint matrix & Loop closure, contact & Coupling and redundancy \\
\hline
Null space $\ker(J_c)$ & Motor abundance & Redundant manipulator & Unactuated DOF \\
\hline
Cotangent space $T^*Q$ & Wrench space & Force and torque & Dual to velocities \\
\hline
Drift field $f$ & Passive dynamics & Gravity, friction & Uncontrolled motion \\
\hline
Input field $g$ & Actuation & Joint torques, muscle forces & Controlled motion \\
\hline
Lie bracket $[g_i, g_j]$ & Controllability coupling & Two-fingered grasp & Motion achievability \\
\hline
Exponential map $\exp$ & Rigid transformation & $T = e^{\hat{\xi}}$ & Pose update \\
\hline
Coadjoint action $\text{Ad}^*$ & Wrench transformation & Frame change in forces & Dual transformation \\
\hline
Momentum map $\mathbf{J}$ & Conserved quantity & Angular momentum & Invariant along trajectories \\
\hline
Geometric phase & Net rotation from shape change & Falling cat, golf swing & Orientation change \\
\hline
\end{tabular}
\caption{Unifying differential geometry and mechanics.}
\label{tab:ch9:mapping}
\end{table}

\subsection{Workflow: From Problem to Solution}

\index{Workflow!analysis and design}%
A typical workflow for analyzing and designing control for a constrained mechanical system:

\subsubsection{Step 1: Problem Formulation}

\begin{enumerate}
  \item \textbf{Identify constraints}. Write down all holonomic constraints $\Phi_j(q) = 0$ (position
  constraints that reduce configuration space). Distinguish from nonholonomic constraints (velocity constraints
  that do not integrate).

  \item \textbf{Verify regularity}. Confirm that the constraint Jacobian $J_c(q)$ has full row rank
  at all feasible configurations (or at least on the region of interest). If not, expect singularities
  and plan accordingly.

  \item \textbf{Compute the constraint manifold dimension}. $n_c = n - \text{rank}(J_c)$, the dimension of
  the configuration space on which the system can move.
\end{enumerate}

\subsubsection{Step 2: Dynamics Construction}

\begin{enumerate}
  \item \textbf{Construct the mass matrix} $M(q)$. Use:
  \begin{itemize}
    \item \textbf{Composite-body algorithm} (Craig 2005): Efficient recursive method for serial chains.
    \item \textbf{Lagrangian formulation}: $M_{ij} = \sum_k (m_k J_{k,i}^\top J_{k,j} + \text{tr}(\cdots))$
    for parallel mechanisms.
    \item \textbf{Numerical differentiation}: Differentiate kinetic energy expressions (slow but general).
  \end{itemize}

  \item \textbf{Compute Coriolis and centrifugal terms}. Either:
  \begin{itemize}
    \item Compute Christoffel symbols from $M(q)$ and use $C_{ij} = -\sum_k \Gamma^i_{jk} \dot{q}_j \dot{q}_k$.
    \item Use $C = \dot{M} - \frac{1}{2} \text{sym}(\text{grad}_q M)$, a shortcut formula.
  \end{itemize}

  \item \textbf{Define gravity and potential energy}. $V(q) = -\sum_k m_k g \, h_k(q)$,
  where $h_k$ is the height of link $k$'s center of mass.

  \item \textbf{Formulate the constrained Euler-Lagrange equations}:
  \begin{equation}
    M(q)\ddot{q} + C(q,\dot{q})\dot{q} + \nabla V(q) = J_c(q)^\top \lambda + B(q)\mathbf{u},
  \end{equation}
  where $\lambda$ are constraint multipliers and $B(q)\mathbf{u}$ is the actuation.
\end{enumerate}

\subsubsection{Step 3: Control Analysis}

\begin{enumerate}
  \item \textbf{Identify actuators and input matrix}. Determine which degrees of freedom are actuated:
  $B(q) \in \mathbb{R}^{n \times m}$.

  \item \textbf{Analyze controllability}. Compute Lie brackets of input fields $g_i = M^{-1} B e_i$ and check:
  \begin{equation}
    \text{rank}\{g_i, [g_i, g_j], [[g_i, g_j], g_k], \ldots\} \stackrel{?}{=} n_c.
  \end{equation}
  If rank is less than $n_c$, the system is uncontrollable; plan differently.

  \item \textbf{Compute accessibility algebra}. The Lie algebra generated by input fields characterizes
  what motions are achievable (important for grasping and manipulation).

  \item \textbf{Analyze null-space redundancy} (if $m > n_c$). Characterize the null space of $B(q)$ or
  the input distribution. Identify which motions can be performed without affecting the main task.
\end{enumerate}

\subsubsection{Step 4: Trajectory Planning}

\begin{enumerate}
  \item \textbf{Define cost function}. Choose:
  \begin{itemize}
    \item Time-optimal: $\int_0^T dt$.
    \item Energy-optimal: $\int_0^T u^\top u \, dt$.
    \item Tracking: $\int_0^T \|q(t) - q_d(t)\|^2 + \|u(t)\|^2 \, dt$.
    \item Biomechanical: Minimize metabolic cost (Zajac 1989) or muscle fatigue.
  \end{itemize}

  \item \textbf{Apply optimal control method}:
  \begin{itemize}
    \item \textbf{Pontryagin Maximum Principle}: Derive necessary conditions for optimality (costate equations).
    \item \textbf{iLQR/DDP}: Iteratively linearize and solve quadratic sub-problems (practical for complex systems).
    \item \textbf{Trajectory optimization}: Direct methods that discretize and solve as finite-dimensional NLP.
  \end{itemize}

  \item \textbf{Project onto constraint manifold}. At each iteration, ensure generated trajectories satisfy
  $\Phi(q(t)) = 0$ and $J_c(q)\dot{q} = 0$ using projection methods or by reparameterizing with reduced coordinates.

  \item \textbf{Validate optimality}. Check costate equations and verify that the Hamiltonian is constant
  along optimal trajectories (for autonomous systems).
\end{enumerate}

\subsubsection{Step 5: Numerical Simulation and Validation}

\begin{enumerate}
  \item \textbf{Choose a constraint-preserving integrator}:
  \begin{itemize}
    \item SHAKE: Fast, simple, but velocity drift over long runs.
    \item RATTLE: More stable, preserves both position and velocity constraints.
    \item Projection method: Exact up to machine precision, but expensive.
  \end{itemize}

  \item \textbf{Integrate forward in time}:
  \begin{equation}
    (q(t_{n+1}), \dot{q}(t_{n+1})) = \text{INTEGRATOR}(q(t_n), \dot{q}(t_n), \Delta t, \text{system parameters}).
  \end{equation}

  \item \textbf{Verify constraint satisfaction}. Check:
  \begin{itemize}
    \item Position: $\|\Phi(q(t))\| < \epsilon_{\text{constraint}}$ (small tolerance, typically $10^{-6}$ to $10^{-8}$).
    \item Velocity: $\|J_c(q(t)) \dot{q}(t)\| < \epsilon_{\text{velocity}}$.
    \item Energy: $\|E(t) - E(0)\| < \epsilon_{\text{energy}}$ (energy should be conserved for passive systems).
  \end{tolerances}

  \item \textbf{Sensitivity analysis}. Vary system parameters (masses, inertias, friction) and re-simulate
  to understand robustness.

  \item \textbf{Compare with experimental data} (if available). For gait analysis or sports biomechanics,
  compare predictions with motion capture and GRF measurements.
\end{enumerate}

\subsubsection{Step 6: Control Implementation}

\begin{enumerate}
  \item \textbf{Synthesize feedback control}. Use the value function or costate trajectory to design
  feedback laws:
  \begin{equation}
    \mathbf{u}(t) = \pi(q(t), \dot{q}(t), t; \theta^*)
  \end{equation}
  where $\theta^*$ are the optimal parameters (e.g., feedback gains from iLQR).

  \item \textbf{Implement on hardware}. Deploy the control law on the actual system (robot or human
  neuromuscular system, via visual feedback and learning).

  \item \textbf{Online adaptation}. If the system is uncertain (e.g., object mass unknown), use
  adaptive or learning-based control to refine the law online.

  \item \textbf{Validate safety}. Ensure the control law respects bounds on forces, torques, and
  accelerations; include fail-safes to prevent damage.
\end{enumerate}

\subsection{Research Directions and Open Problems}

\index{Open problems}%
Several frontier topics extend the theory developed here:

\begin{itemize}
  \item \textbf{Constraint Manifold Learning}. Many systems have implicit constraints (e.g., contact
  constraints with soft materials). Learning the constraint manifold from data, and updating control
  design accordingly, is an active area (Kakade et al., and the broader learning dynamical systems community).

  \item \textbf{Stochastic Geometric Mechanics}. How do noise and uncertainty propagate on a constraint
  manifold? Stochastic differential geometry provides tools, but their application to control is nascent.

  \item \textbf{Motor Learning as Riemannian Optimization}. Can human motor learning be understood as
  gradient descent on the constraint manifold with respect to some cost (energy, accuracy, smoothness)?
  Preliminary work suggests yes (Graziano, Aflalo, and neuromotor control literature).

  \item \textbf{Geometric Phase Optimization for Legged Locomotion}. Rather than minimizing center-of-mass
  work, can we design gaits that optimize the geometric phase to reduce torque requirements or improve stability?

  \item \textbf{Singularity Avoidance}. Real-time planning that explicitly avoids configurations where
  the constraint Jacobian loses rank, maintaining controllability margin throughout trajectories.
\end{itemize}

\section{Summary and Further Reading}

This chapter has developed the geometric framework for constrained mechanical systems, connecting differential
geometry, mechanics, control theory, and numerical methods. The key insights are:

\subsection{Core Concepts and Structures}

\begin{itemize}
  \item \textbf{Configuration manifolds and holonomic constraints} define the feasible motion space. A constraint
  manifold $Q_c = \{\Phi(q) = 0\}$ is a lower-dimensional submanifold of the unconstrained space, with dimension
  $\dim(Q_c) = n - \text{rank}(J_c)$. Regular constraints ensure smooth submanifold structure.

  \item \textbf{Lagrangian mechanics on submanifolds} encodes the cost of motion via kinetic energy. The mass matrix
  $M(q)$ induces a Riemannian metric, and the Christoffel symbols capture how basis vectors rotate on the curved
  configuration space. These geometric effects (Coriolis, centrifugal) arise naturally from the manifold curvature.

  \item \textbf{Control-affine systems with drift and input fields} ($\ddot{q} = f(q,\dot{q}) + \sum_i g_i u_i$)
  allow systematic analysis of controllability via Lie brackets. The Lie algebra generated by input fields determines
  what motions are achievable.

  \item \textbf{Null spaces and redundancy} structure enable motion planning in high-dimensional systems. When actuators
  exceed degrees of freedom ($m > n_c$), the null space $(I - B^\dagger B)$ parameterizes motions that preserve the
  primary task. Null-space optimization allows secondary objectives (energy minimization, singularity avoidance).

  \item \textbf{Screw theory and Lie groups} provide a coordinate-free language. Twists (infinitesimal rigid motions)
  live in $\mathfrak{se}(3)$, wrenches in the dual $\mathfrak{se}^*(3)$. The adjoint and coadjoint actions transform
  these objects under reference frame changes, enabling force analysis in multiple frames.

  \item \textbf{Optimal control principles (Pontryagin Maximum Principle)} guide trajectory design. The Hamiltonian
  encodes the trade-off between cost (running cost + control effort) and dynamics (drift and actuation). Solving the
  resulting costate ODEs yields the optimal control law and value function.

  \item \textbf{Observability and identifiability} characterize what can be learned from measurements. The observability
  rank condition (via Lie derivatives) determines which state components are uniquely determined by output trajectories.
  The unobservable subspace creates ambiguity; resolving this requires additional measurements (EMG, optimization criteria).

  \item \textbf{Connections and geometric phase} reveal subtle rotational effects from shape changes. Even in the absence
  of external torques, coordinated shape changes (e.g., arm raises in a golf swing) accumulate net rotation via the
  mechanical connection. The accumulated phase is proportional to the area swept in shape space.

  \item \textbf{Numerical methods (SHAKE, RATTLE, SVD projection)} make theory actionable. Constraint-preserving integrators
  enforce constraints up to machine precision while maintaining energy conservation. SVD provides stable null-space projection.
\end{itemize}

\subsection{Unified Perspective}

All of the above structures are manifestations of a single geometric principle: \emph{motion on a manifold is constrained
by its intrinsic geometry}. The metric $M(q)$ encodes what directions are ``expensive'' to move; the constraint Jacobian
$J_c(q)$ encodes feasible directions; Lie brackets encode achievable changes via control. Optimal control solves for the
trajectory that minimizes cost while respecting these geometric constraints.

\subsection{Applications}

The framework applies to diverse problems:
\begin{itemize}
  \item \textbf{Parallel mechanisms and grippers}: Understand force/velocity trade-offs, singularities, and grasp stability.
  \item \textbf{Legged locomotion}: Analyze gait dynamics under bilateral foot contact; optimize for energy efficiency.
  \item \textbf{Sports biomechanics}: Predict muscle activations, understand geometric effects (e.g., rotation via geometric phase).
  \item \textbf{Humanoid robotics}: Plan whole-body motion respecting balance and manipulation constraints simultaneously.
\end{itemize}

\subsection{Further Reading}

For deeper study, key references are:

\begin{enumerate}
  \item \cite{BulloLewis2004}: \textit{Geometric Control of Mechanical Systems}. Comprehensive treatment of Lie groups,
  differential geometry, and optimal control for mechanics. Essential for understanding mechanical connections and
  geometric phase.

  \item \cite{AbrahamMarsden1978}: \textit{Foundations of Mechanics}. Foundational text on differential geometry and
  Lagrangian mechanics. Rigorous mathematical treatment of manifolds, tangent spaces, and Lie groups.

  \item \cite{Arnold1989}: \textit{Mathematical Methods of Classical Mechanics}. Elegant geometric perspective on classical
  mechanics. Covers configuration manifolds, symplectic geometry, and conservation laws via Noether's theorem.

  \item \cite{Murray1994}: \textit{A Mathematical Introduction to Robotic Manipulation}. Practical guide to kinematics,
  dynamics, and control of manipulators. Screw theory and Jacobian analysis.

  \item \cite{marsden1999}: \textit{Introduction to Mechanics and Symmetry}. Focuses on symmetries and conservation laws;
  includes discussion of Routh reduction and geometric phase.

  \item \cite{pontryagin1962}: \textit{The Mathematical Theory of Optimal Processes}. Classic monograph on Pontryagin
  Maximum Principle and optimal control. Foundational for trajectory optimization.

  \item \cite{HairerLubichWanner2006}: \textit{Geometric Numerical Integration}. Comprehensive treatment of structure-preserving
  integrators. Covers symplectic methods, constraint-preserving algorithms (SHAKE, RATTLE), and error analysis.

  \item \cite{Zajac1989}: \textit{Muscle and Tendon: Properties, Models, Scaling, and Applications to Biomechanics and
  Motor Control}. Comprehensive biomechanics reference. Details muscle physiology, force-length-velocity relationships,
  and optimization methods for inverse dynamics.

  \item \cite{TodorovJordan2002}: \textit{Optimality Principles in Sensorimotor Control}. Elegant formulation of motor
  control as optimal control under uncertainty. Discusses how the nervous system may implement near-optimal strategies.

  \item \cite{deLeva1996}: \textit{Adjustments to Zatsiorsky-Seluyanov's Segment Inertia Parameters}. Provides anthropomorphic
  data (segment masses, inertias) for human biomechanics simulations.

  \item \cite{ball1900}: \textit{Theory of Screws}. Classic treatise on screw theory in rigid body mechanics. Foundational
  for modern treatments of twists and wrenches.

  \item \cite{Goldstein2002}: \textit{Classical Mechanics} (3rd ed.). Comprehensive classical mechanics text. Covers Lagrangian
  and Hamiltonian formalism, rigid body dynamics, and introduction to differential geometry.
\end{enumerate}

\subsection{Looking Forward}

The geometric methods of this chapter—configuration manifolds, Riemannian metrics, optimal control—are increasingly
recognized as the right language for complex, multi-body systems. As systems grow more complex (humanoid robots with
many joints, biological systems with hundreds of muscles), geometric structure provides the essential simplifications
that make analysis and design feasible. Future work will increasingly combine these classical geometric methods with
modern machine learning, using data to learn constraint manifolds and cost functions automatically.